\documentclass[../main.tex]{subfiles}
\begin{document}
%==========================================TIÊU ĐỀ TẬP========================================
\begin{table}[h]
  \begin{adjustwidth}{-1cm}{}
    \begin{tabular}{>{\centering\arraybackslash}p{6cm}|>{\raggedright\arraybackslash}p{12.5cm}}
      \multirow{5}{6cm}
      % Cột bên trái: ÉPISODE và số tập
      {\centering
        {\\[-40pt]
          \flaregothic\fontsize{20pt}{20pt}\selectfont
          \color{eptype}ÉPISODE\color{black}\\
          \burbank\fontsize{72pt}{72pt}\selectfont
          \color{epnum}93
          \\[8pt]
          \hrule
          \vspace{8pt}
          \flaregothic\fontsize{20pt}{20pt}\selectfont
          \color{eptype}SPÉCIAL\color{black}\\
          \burbank\fontsize{72pt}{72pt}\selectfont
          \color{epnum}XVI
          \\[18pt]
          % \flaregothic\fontsize{14pt}{14pt}\selectfont
          % \color{epattr}BIENVENUE, \uline{\color{Banpire} Mel} !
          \flaregothic\fontsize{18pt}{18pt}\selectfont
          \color{holop}L'HOLOPROJET !\\
          \color{black}
        }
      }
       &
      % Dòng thứ nhất: tên tập tiếng Nhật
      {\fotskip\fontsize{18pt}{18pt}\selectfont \ruby{助}{たす}けて！クビにされる！
      }                       \\[6pt]
      % Dòng thứ hai: tên tập tiếng Anh
       & {\cafeteria\fontsize{18pt}{18pt}\selectfont "Bloom," - Declaration Misleading}     \\[4pt]
      % Dòng thứ ba: tên tập tiếng Việt
       & {\vnmsans\fontsize{18pt}{18pt}\selectfont Bản kế hoạch}                            \\[8pt]
      % Dòng thứ tư: Nhân vật xuất hiện
       & {\brian{\fontsize{14pt}{14pt}\selectfont Track used: Neon Mixtape Tour - Zombotany}
         }
      \\[8pt]
      % Dòng cuối cùng: Thời gian diễn ra của tập (thường sẽ là ngày phát hành)
       & {\continuum\fontsize{16pt}{16pt}\selectfont Ngày phát hành: 14/02/2021}            \\[36pt]
    \end{tabular}
  \end{adjustwidth}
\end{table}
%=========================================TRUYỆN CHÍNH========================================
\color{black}

Để chuẩn bị cho buổi live concert online \textit{Bloom,} - một sự kiện đặc biệt của \hll\, nơi các thành viên có thể phô diễn tài năng của mình - ai nấy đều đang gấp rút luyện tập. Nhưng giống như mọi lần, không khí chuẩn bị chẳng bao giờ có thể yên bình như mong đợi.

Sự bát nháo trong nội bộ trở thành điều không thể tránh khỏi. Người thì quên lịch tập, kẻ thì chạy khắp nơi tìm trang phục biểu diễn, có người còn lỡ tay làm đổ nước lên thiết bị thu âm khiến bộ phận kỹ thuật suýt phát khóc. Trong khi đó, một số thành viên lại mải mê chỉnh sửa lời bài hát hoặc tập vũ đạo đến mức quên cả ăn uống, để rồi đến khi kiệt sức thì mới nhận ra sai lầm.

Mio, Flare, Matsuri, Kanata, Aqua và Marine cũng như vậy.

Và điều đáng nói nhất là\dots\ họ đang không có một người quản lý thay thế.

Hikawa Kuon - quản lý phụ trách điều hành chung - đã xin phép nghỉ về quê ăn Tết Nguyên Đán từ một tuần trước, ngay sau khi đã giải quyết hậu quả từ Botan, Luna, Towa và chính Marine. Quyết định này không hề nằm ngoài dự đoán, bởi cậu đã làm điều tương tự vào năm ngoái, và cũng như lần trước, YAGOO đã dễ dàng chấp thuận.

Nhưng điều không ai ngờ tới là, chính vì quyết định này mà cậu đã vô tình chạm trán với một người mà đáng ra cậu không thể gặp - một người thuộc về thế hệ đã từng có năm thành viên nhưng giờ đây chỉ còn bốn.

Cuộc gặp gỡ đó sẽ còn thay đổi rất nhiều thứ, nhưng\dots\ đó lại là một câu chuyện khác, sẽ được kể ở **Quyển D: Ký túc xá nhà Hikawa**.

Còn hiện tại, câu chuyện vẫn tiếp tục ở đây - nơi các thành viên \hll\ đang đối mặt với một vấn đề hoàn toàn khác. Một vấn đề không liên quan đến những buổi luyện tập căng thẳng hay sự vắng mặt của quản lý, mà là một tờ giấy kỳ lạ\dots\ Một tờ giấy mà chẳng ai biết nó đến từ đâu, nhưng chắc chắn rằng, nội dung trên đó có khả năng làm chấn động cả công ty.

\il

Một tiếng nổ vang rền khắp văn phòng, ngay sau đó là hàng loạt tiếng hét thất thanh vang lên từ nhiều phía. Người thì chạy tán loạn, người thì ngã sõng soài xuống sàn nhà, đồ đạc bị hất tung tóe khắp nơi.

Một người nào đó thốt lên trong bất lực: ``\textbf{Mọi người à!! Bình tĩnh lại được không!?}''

Nhưng làm gì có ai còn giữ được bình tĩnh trong cái mớ hỗn độn này nữa.

Thật ra, đây cũng chẳng phải là lần đầu tiên văn phòng \hll\ rơi vào cảnh náo loạn như vậy, nhất là trong thời gian chuẩn bị cho một sự kiện quan trọng như \textit{Bloom,}. Với việc mọi người tất bật luyện tập, xử lý trang phục, chuẩn bị sân khấu, còn ban quản lý thì chạy đôn chạy đáo lo hậu cần, hỗn loạn gần như đã trở thành một phần không thể thiếu của công ty này.

Nhưng lần này thì có lẽ mọi chuyện đã đi hơi quá. Không ai rõ chuyện gì vừa mới xảy ra. Chỉ biết rằng, một vụ nổ đã xảy ra, mọi thứ thì rối tung lên, và giờ họ phải bắt tay vào dọn dẹp trước khi tình hình trở nên tệ hơn.

\ct

``Lại một ngày hỗn loạn nữa rồi\dots'' Flare thở dài trong khi đang giữ cây chổi lau nhà, mắt nhìn đồ đạc bị rơi vãi khắp nơi.

Mio đứng gần đó, đang quét dọn với vẻ mặt đầy cam chịu. Cô ho một tiếng do hít phải bụi rồi cũng than thở: ``Tại sao mọi chuyện lại thành ra thế này chứ?''

Matsuri thì ôm một hộp bìa khá lớn, lảo đảo đi về phía hai người và lẩm bẩm: ``Chúng ta đáng ra phải là idol chứ, nhưng nhìn xem\dots sao giống mấy nàng hề hơn vậy\dots''

\img{7-11a}

Nàng Yêu tinh vàng bật cười nhẹ, vừa quét sàn vừa đáp: ``Thì, \hll\ là như vậy mà.''

Nàng Sói đen cũng gật đầu, tặc lưỡi: ``Giống y như những gì Kuon-senpain nói.''

Kanata đứng kế bên, đang bê một hộp khác, bật cười trước câu nói đó. Nhưng chưa kịp góp chuyện thì Mio đã nhắc nhở: ``Tập trung làm việc đi, kẻo đến muộn thì còn lâu mới dọn xong đống này.''

Cả hai người đàn em của cô đồng thanh đáp với tông giọng lười biếng: ``Rồi\~{}''

Matsuri đang định tiếp tục công việc thì chợt nhận thấy có hai tờ giấy rơi ra từ đống hộp mà cô đang xếp.

\img{7-11b}

Cô cúi xuống nhặt lên, nhìn chằm chằm vào chúng một lúc lâu, rồi gọi ba người kia lại: ``Nè nè, mấy người có biết đây là cái gì không?''

Cô giơ tờ giấy ra, trên mặt có vẻ khó hiểu: ``Chị tìm thấy nó khi đang lau dọn này.''

Cả ba người còn lại cũng tò mò tiến lại gần để xem xét. Mio cầm lấy một trong hai tờ giấy, nheo mắt nhìn kỹ hơn, rồi chợt thốt lên: ``Xem nào, xem nào\dots\ À.''

``Sao thế?''

Nàng Sói đen không trả lời ngay, chỉ im lặng một lúc rồi nghiêng đầu đáp lại: ``Cái này là của Kuon-senpai đây mà.''

``Ể? Của anh ấy sao?'' Kanata ngạc nhiên.

Mio gật đầu, giơ tờ giấy lên để mọi người nhìn rõ hơn: ``Đây là đơn nghỉ phép của anh ấy.''

Không khí xung quanh thoáng im lặng một chút.

``Nhưng sao đơn nghỉ phép của anh ấy lại ở đây?'' nàng Thiên sứ nhíu mày, nhìn chằm chằm vào tờ giấy.

Câu hỏi đó làm tất cả bốn người đều cảm thấy có chút khó hiểu. Dù thế, chỉ có một người trong số bốn cô gái lại biết được lý do đằng sau.

``Chuyện là như thế này\dots'' Mio thở dài, khoanh tay lại rồi bắt đầu kể vắn tắt câu chuyện.

\begin{center}
  \uline{Hồi tưởng\\Một tuần trước.}
\end{center}

Sau khi giải quyết xong vụ rắc rối giữa Botan, Towa và Marine, Kuon thở hắt ra, lấy lại bình tĩnh rồi đi thẳng đến phòng giám đốc - nơi YAGOO đang làm việc. Mục đích của cậu rất rõ ràng: xin phép nghỉ Tết Nguyên Đán để về quê.

Tất nhiên, cậu đã dự tính trước mọi khả năng. Nhưng giống như năm ngoái, YAGOO nhanh chóng đồng ý, không hề do dự. Ông hiểu rõ truyền thống của gia đình Kuon, và cũng biết rằng cậu chưa bao giờ làm trái những quy tắc ấy.

Sau khi nhận được sự chấp thuận, Kuon quay lại văn phòng, trên tay cầm hai tờ giấy - một là bản gốc mang về, một là bản sao cậu để lại. ``Anh để lại bản sao tờ xin nghỉ phép của anh nha.''

Mio, người đang sắp xếp giấy tờ ở bàn làm việc, liếc nhìn cậu, chau mày: ``Anh có nhất thiết phải làm thế không?''

Cậu Tổng nhún vai, giọng điệu thoải mái: ``Cần chứ. Để cho mấy đứa khác vào mà nếu thắc mắc anh ở đâu thì cứ giơ tờ nghỉ phép đó cho mọi người xem.''

Mio khẽ thở dài. ``Em thấy nó thừa thãi thế nào ấy\dots\ Em có thể nói một tiếng với họ là được mà.''

``Thừa còn hơn thiếu,'' cậu cười nhẹ, đặt tờ giấy xuống một góc bàn rồi vỗ nhẹ lên đó.

``Thôi thì, em sẽ giữ nó vậy.'' Nàng Sói đen chống cằm, suy nghĩ một chút rồi hỏi: ``Vậy khi nào anh về?''

Cậu Tổng ngẫm nghĩ một hồi như đang tính toán, rồi đáp: ``Tầm mùng 8, mùng 9 Tết. Hôm nay là 26 Âm lịch, vậy là khoảng\dots''

Cô nàng nhẩm tính trong đầu rồi nhanh chóng kết luận: ``12 đến 13 ngày. Tức là sau buổi concert của bọn em à?''

Kuon gãi đầu, vẻ hơi áy náy: ``Ừ\dots\ Xin lỗi mấy đứa nha. Anh cũng muốn giúp mấy đứa vào buổi live sắp tới, nhưng ở nhà anh, gia đình luôn ưu tiên hơn so với công việc. Giám đốc cho anh nghỉ về quê ăn Tết đã là đãi ngộ rồi.''

Cậu cười nhẹ, nhìn Mio đầy tin tưởng: ``Anh nghĩ là mấy đứa sẽ đàng hoàng lúc anh về thôi mà, không phải lo cho anh đâu.''

Hôm đó, mặc dù trời mưa nặng hạt, Kuon vẫn nhanh chóng thu xếp hành lý, vội vã bắt xe buýt về thành phố Kawachi.

\begin{center}
  \uline{Kết thúc hồi tưởng.}
\end{center}

Flare chậm rãi gật đầu, giọng có phần hiểu ra vấn đề: ``Hóa ra là vậy\dots''

Matsuri thả mình xuống ghế, chống tay lên cằm, tặc lưỡi: ``Văn hóa ăn Tết Âm lịch nhà anh ấy là vậy mà.''

Kanata thở dài, lắc đầu: ``Chán thật nhỉ.''

Bốn người ngồi im lặng một chút, rồi lại nhìn tờ đơn nghỉ phép của Kuon, cảm thấy có chút hụt hẫng khi người quan trọng nhất - chỉ sau YAGOO - lại không có mặt ở đây.

``\dots Mà, còn một tờ giấy nữa này,'' nàng Thiên sứ chỉ vào tờ giấy còn lại mà Sói đen đang cầm.

Matsuri giở tờ giấy nằm sau tờ đơn ra. Trên đó ghi tiêu đề: ``Báo cáo và danh sách thay thế về việc mở rộng và ảnh hưởng trong năm nay''.

Sói đen nhíu mày khó hiểu: ``Ủa? Gì đây?''

Họ cùng nhìn xuống nội dung, nhưng nhiều chỗ đã bị bôi đen hoặc làm mờ: ``Mở rộng [REDACTED] quá nhanh lên đến 100 [REDACTED] đã tạo ra ảnh hưởng bất lợi\dots''

Nàng Dị giới suy đoán: ``Cái này nói về chúng ta sao?''

Vừa dứt lời, Kanata mở to mắt. Flare chỉ vào một đoạn trong báo cáo: ``Đây này, con số 100 này có thể nói đến việc vượt quá 100 triệu lượt xem, kiểu như vậy.''

\img{7-11c}

Mio tiếp tục đọc, sắc mặt cô dần trở nên nghiêm trọng: ``[REDACTED] và [REDACTED] thì chỉ có cười, điều đó cho thấy chúng ta đã đi lệch quá nhiều so với mục tiêu ban đầu.''

``Cái này nghe giống như đang than phiền rằng chúng ta cứ làm trò hề quá không?''

Nàng Thiên sứ bất giác thấy lo lắng: ``\dots Vậy cái phân đoạn cần loại bỏ là thế nào chứ?''

Cô chợt giật mình khi nhìn thấy một câu trong báo cáo: ``Do đã được xếp hạng D của [REDACTED], chúng ta nên loại bỏ [REDACTED].''

Thiên sứ hét toáng lên: ``Không lẽ là\dots\ đuổi việc!?''

Mio giật bắn mình, cả toàn thân cô đông cứng lại. Matsuri vội trấn an: ``Không đời nào đâu! Chị đã là một idol chính hiệu rồi mà!''

Nhưng Flare lại nói một câu khiến mọi hy vọng sụp đổ: ``Nhưng trong này có ghi không thể để thế này được' còn gì.''

Matsuri lập tức câm nín. Trong một khoảnh khắc, thế giới xung quanh họ như sụp đổ. Cả ba người đồng loạt tuyệt vọng thốt lên những điều trăn trối mà họ chưa làm được:

\begin{dialogue}
  \speak{Mio} *ôm đầu thất thần, la hét trong thảm thiết* \textbf{Waaaah!!}
  \speak{Matsuri} *khuỵu xuống* \textbf{Chị vẫn chưa được nhìn quần lót của Lamy-chan mà!}
  \speak{Kanata} *ôm đầu thất vọng* \textbf{Em còn chưa trả hết nợ tiền PC nữa!}
\end{dialogue}

Riêng nàng Dị giới vẫn đứng đó, vẻ mặt tỏ ra khó hiểu trước sự hoảng loạn của ba người.

\img{7-11d}

Cô liền lên tiếng trấn an: ``Chắc phải có sự nhầm lẫn gì chứ\dots\ Sao chúng ta không thử hỏi Kuon-senpai xem? Anh ấy là Tổng quản lý chúng ta mà, chắc chắn phải biết chuyện này!''

Nhưng Sói đen chỉ lắc đầu, giọng nói đầy tuyệt vọng: ``Vô ích thôi. Dù có gọi thế nào thì anh ấy cũng sẽ không nhấc máy đâu.''

Flare nhíu mày: ``Sao lại thế?''

Nàng Sói đen thở dài, giọng chua xót: ``Giờ anh ấy đang nghỉ Tết. Anh ấy để máy liên lạc với chúng ta ở nhà rồi, chỉ có thể gọi điện và trao đổi với A-chan hoặc YAGOO thôi.''

``Chính chị đã thử gọi cho anh ấy mấy ngày trước đó nhưng cuối cùng anh ấy có nhấc máy đâu\dots'' nàng Cổ động ở bên cạnh nói thêm, đầu vẫn rụp xuống.

Flare sững sờ: ``Điều đó có nghĩa là\dots''

Mio gật đầu, chốt lại: ``Dù có gọi hay nhắn tin thì anh ấy cũng không trả lời đâu.''

Không tin vào điều đó, nàng Dị giới thử nháy máy gọi cho cậu, nhưng chỉ nhận được một câu trả lời vô cảm từ tổng đài: ``Thuê bao quý khách vừa gọi hiện không liên lạc được. Xin quý khách vui lòng gọi lại sau.''

Cô thử nhắn tin, nhưng chỉ thấy hiện trạng thái ``đã gửi'' mà không phải ``đã xem''.

Lúc này, cô hoàn toàn hiểu ra vấn đề. Cô đứng lặng trong cơn sốc, nhìn ba người kia kêu la trong tuyệt vọng. Cô bất giác thở dài: ``Mọi người có thể bình tĩnh lại một chút được không?''

Ngay lúc này, Matsuri bật dậy, thốt lên: ``Bình tĩnh làm sao được chứ!? Đây có thể là ngày cuối cùng của chúng ta đấy!!''

Mio thì vẫn đang ôm đầu, ỉ ôi: ``Chuyện gì đang xảy ra vậy trời\dots''

Kanata thì nước mát lưng tròng: ``Kuon-senpai bỏ rơi chúng ta thật rồi\dots''

Nàng Dị giới bối rối trước tình cảnh trước mắt. Một bản báo cáo bí ẩn, một tờ đơn nghỉ phép không rõ tại sao lại xuất hiện ở đây, và một nhóm idol đang hoảng loạn. Không lẽ mọi chuyện lại nghiêm trọng đến thế sao?

Flare lật lại tờ giấy, cố gắng tìm thêm thông tin. Nhưng những phần bị che mất khiến cô khó có thể hiểu rõ toàn bộ nội dung.

``Này, các cậu có chắc là chúng ta không đang hiểu lầm gì không? Có khi nào đây chỉ là một tài liệu nội bộ chưa hoàn chỉnh thôi?'' cô lên tiếng, nêu ra giả thuyết.

Nàng Cổ động ngước mắt lên nhìn vào cô: ``Cậu nghĩ vậy thật à?''

``\dots Ừ thì\dots'' cô đổ mồ hôi hột, giải thích có lý nhất có thể: ``ít nhất cũng nên xác nhận lại trước khi hoảng loạn chứ.''

Nàng Sói đen thở dài: ``Nhưng chúng ta đâu còn cách nào để hỏi nữa\dots''

Không ai trong số họ có thể liên lạc với Kuon. Nếu muốn biết sự thật, họ cần một người khác có thể giải thích về tài liệu này.

Bất chợt, Kanata chợt lóe lên một ý tưởng: ``A-chan! Chúng ta có thể hỏi A-chan mà!''

Nghe vậy, đôi mắt của nàng Cổ động rực sáng, thắp lên niềm hy vọng: ``Đúng rồi! Chị ấy là trợ lý giám đốc mà! Nếu ai biết chuyện này thì chắc chắn là chị ấy!''

Mio gật đầu: ``Vậy thì còn chần chừ gì nữa? Đi thôi!''

Bốn người lập tức thu dọn giấy tờ, vội vã rời khỏi phòng để tìm A-chan. Họ chạy dọc hành lang, hơi thở gấp gáp vì căng thẳng.

Khi đến cửa phòng làm việc cá nhân của cô, Kanata vội vàng mở cửa mà không thèm gõ trước. ``\textbf{A-chan! Chị có ở đây không!?}''

Nhưng căn phòng trống trơn. Không có bóng dáng của cô trợ lý quen thuộc, chỉ còn lại chiếc bàn làm việc gọn gàng với một vài tài liệu xếp ngay ngắn. Đèn trong phòng cũng đã tắt từ lâu, cho thấy chủ nhân của nó đã không còn ở đây từ trước đó.

``Ủa? A-chan đâu rồi!?'' Matsuri thảng thốt, ngó quanh khắp căn phòng.

Mio chợt sững lại, ký ức chợt ùa về. Cô vỗ mạnh vào trán, vẻ mặt như vừa nhận ra điều gì đó khá quan trọng: ``Chết thật\dots\ Chị ấy được công ty cho nghỉ mấy ngày để dưỡng sức mà!''

``Hả!?'' tất cả mọi người đồng thanh, ngỡ ngàng với câu nói này.

Nàng Sói đen thở dài, giải thích: ``Mấy hôm trước, A-chan làm việc xuyên đêm nhiều quá, suýt ngất luôn trong văn phòng. YAGOO bắt chị ấy nghỉ ngơi vài ngày cho lại sức rồi!''

``Nghĩa là\dots\ bọn mình không thể hỏi chị ấy bây giờ rồi,'' đôi vai của nàng Thiên sứ rũ xuống, bày tỏ sự tuyệt vọng.

Nàng Lễ hội thả người xuống ghế, rầu rĩ: ``Vậy thì phải làm sao đây? Bọn mình không hiểu hết nội dung báo cáo, cũng không liên lạc được với Kuon-senpai\dots''

``Chắc phải tự luận ra ý nghĩa của nó thôi,'' Mio nói, không giấu nổi sự thất vọng.

``Nhưng có quá nhiều chỗ bị bôi đen thế này\dots'' cô cầm tờ giấy tài liệu nội bộ đó, liếc nhìn những chỗ bị che mất bởi mực đen.

Cả bốn người nhìn nhau đầy lo lắng. Họ chẳng có đầu mối nào cả, mà càng suy nghĩ, họ lại càng cảm thấy bất an.

Ngay lúc đó\dots

Cạch.

Tiếng cửa mở vang lên, thu hút sự chú ý của cả bọn. Họ đồng loạt quay đầu lại, chỉ để thấy hai bóng người rình rập ở bên ngoài.

Marine và Aqua ló mặt vào, chưa hiểu chuyện gì đã xảy ra với bọn họ. Nàng Thuyền trưởng lên tiếng trước: ``Ủa? Mọi người làm gì mà mặt mày thất thần vậy?''

Nàng Hầu gái tiếp lời sau đó: ``Sao trông ai cũng như mất hết ý chí sống thế này?''

Cả bốn người trong phòng đồng loạt quay ra nhìn hai kẻ vừa mới đến bằng ánh mắt vô hồn, trống rỗng như những con người đã đánh mất tất cả. Không gian trầm lắng đến mức khiến Marine bất giác rùng mình: ``C-Cái gì vậy!? Đừng nhìn em bằng ánh mắt đó chứ!''

\img{7-11e}

Không ai trong nhóm trả lời. Họ chỉ lặng lẽ đưa cho Marine và Aqua tờ báo cáo đầy rẫy những chỗ bị bôi đen, đồng thời thều thào bằng giọng tuyệt vọng: ``Chúng ta\dots\ sắp bị đuổi việc rồi\dots''

``\textbf{Hả!!?}''

Cả hai người vội vã cầm lấy tờ giấy, đọc lướt từ trên xuống dưới. Khi hai người đã hiểu ra phần nào nội dung, họ đồng thanh hét lên: ``\textbf{Cái gì cơ!? Bị đuổi việc á!?}''

``Nhưng mà nhưng mà-!'' Senchou lắp bắp, mắt trố ra không tin vào những gì cô đang đọc, ``em vốn dĩ đã là một idol thanh khiết và đáng mến từ trước tới nay mà!''

Lập tức, nàng Hầu gái nhìn chằm chằm với ánh mắt nghi ngờ: ``Em nói thật á?'' Những người còn lại cũng ngước mắt nhìn cô với những vẻ mặt khó tin.

\img{7-11f}

Sau câu nói đó, Marine đứng hình ngay tại chỗ. Cô cố tìm lời biện hộ, nhưng cuối cùng lại chẳng nói được gì, chỉ quay mặt đi ho khan một tiếng để lảng tránh.

Matsuri thở dài một hơi, vỗ tay lấy lại tinh thần. Cô khoanh tay, nhìn mọi người với vẻ mặt quyết tâm: ``Chị biết là mọi người vẫn còn đang hoảng loạn, nhưng mà bây giờ vẫn chưa quá muộn đâu! Nếu chúng ta chưa đạt tiêu chuẩn idol, thì chỉ còn một cách\dots''

``\dots Cách gì?'' tất cả mọi người đồng thanh hỏi.

``Chúng ta phải luyện tập để cải thiện khả năng làm idol của mình thôi!''

Cả nhóm ngẩn người trước lời tuyên bố đầy quyết tâm ấy, nhưng liệu kế hoạch này có thực sự giúp họ thoát khỏi nguy cơ bị loại bỏ hay không?

\ct

Matsuri, giờ là ``huấn luyện viên'' tự xưng, mang về năm chiếc túi màu vàng, bên trong có chứa đồ để tập luyện.

Mio nhìn đống túi lạ lẫm kia mà chẳng hiểu Matsuri định làm gì, liền hỏi: ``Ê-tô\dots\ Vụ này là sao đây?'

\img{7-11g}

Nàng Cổ động đắc ý đáp lại: ``Bọn chị đã ra ngoài và mua một số mặt hàng kinh điển của idol đây!''

Phía sau cô, bộ ba Aqua, Flare và Kanata đang túm tụm lại trò chuyện.

Nàng Hầu gái ngó vào một trong những chiếc túi, ngờ ngợ hỏi: ``Ở đây có gì thế nhỉ?''

Nàng Dị giới nhún vai: ``Đồ để luyện tập làm idol ấy mà. Em với cả Kanatan đã lựa mấy món phù hợp với bọn mình rồi.''

Nàng Thiên sứ gật gù: ``Khó chọn lắm chứ chẳng đùa. Thôi cứ để đó xem Matsuri-senpai định làm gì với chúng nào.''

Trong khi cả nhóm còn đang tò mò, Marine đứng cạnh, khoanh tay lại với vẻ (có chút) ngạo mạn, mạnh dạn tuyên bố: ``Mặc vào đi đã, rồi hỏi sau!''

Cô còn giơ ngón tay cái lên như thể đang nói điều hiển nhiên nhất thế giới. Mio nhìn Marine, khẽ nhíu mày: ``Sao chị nghe em nói mà thấy đáng ngờ thế nhỉ?''

Trong đầu Sói đen thoáng hiện lên hình ảnh Marine ``chuẩn hải tặc,'' trông chẳng khác gì một kẻ đang chực chờ làm trò lố.

Matsuri vỗ tay một cái, rồi hất tóc đầy tự tin: ``Mio-chan cứ khéo lo! Bọn em đã tự chọn ra những món đồ phù hợp với bản thân rồi!''

Cô chống tay lên hông, quay sang Mio và chỉ tay với vẻ đắc thắng: ``Mio-chan, em sẽ là ban giám khảo nhé!''

Đôi tai sói của cô cụp xuống, cô khoanh tay lại, ánh mắt đầy hoài nghi: ``Không được! Tôi cá là mấy người lại định giỡn nhây với nhau thôi!''

Dù thế, nàng Lễ hội cũng chỉ vẫy tay xua đi: ``Cứ tin tưởng vào bọn này đi\~{}''

Marine hào hứng tiến lên, đặt tay lên ngực mình tuyên bố: ``Vậy để Senchou đây bắt đầu trước ha!''

Mio chỉ còn biết thở dài, ngán ngẩm suy nghĩ: ``\textit{Không biết liệu mọi người có làm được gì ra hồn không nữa\dots}''

\ct

Buổi tập luyện mở màn với cảnh năm người\dots\ trượt giày patin (pa-tanh).

Mio chỉ biết ngồi chống cằm quan sát, lòng thầm nghĩ: ``\textit{Đừng bảo đây là bài kiểm tra idol nhé\dots}''

Tuy nhiên, bọn họ lại vô cùng nghiêm túc (tất nhiên, theo cách rất riêng của họ). Aqua và Kanata lảo đảo bám lấy nhau để giữ thăng bằng, Flare thì cố gắng đứng vững nhưng vẫn lắc lư, còn Marine và Matsuri thì\dots\ dường như đã quen với trò này từ trước.

Với mỗi người trượt về phía trước, họ cố gắng vòng ra đằng sau, để cho Marine - người bày ra trò này - kiễng người xuống, đưa tay lên trên cao, giả vờ ngân nốt cao.

``Đúng rồi đó! Idol chân chính phải có bước di chuyển duyên dáng!'' nàng ``huấn luyện viên'' nhìn cách mọi người đồng điệu với nhau dù còn khá vụng về.

Trong khi đó, nàng Sói đen - người đeo sẵn một cặp kính râm và quàng một chiếc khăn hồng chéo người - ngán ngẩm: ``Tôi không nhớ có điều đó trong giáo trình idol đâu\dots''

Sau vài phút cố gắng, cả nhóm quyết định thử biểu diễn một màn trình diễn nhỏ. Màn trình diễn kết thúc, với năm người họ đứng một cách chỉnh chu.

\img{7-11h}

Và, câu đầu tiên mà ``giám khảo'' Sói đen thốt lên lại là\dots\ ``Tại sao đến cả ý tưởng idol của mấy người lại lạc hậu thế!?'' với một vẻ mặt không hề biến sắc. ``Trông chả khác nào như ở thời Chiêu Hòa cả!''

Nàng Thuyền trưởng thấy vậy, liền tiếp tục động tác trượt patin, giải thích: ``Văn hóa nhạc pop luôn luôn được làm mới! Luôn như vậy-''

Cô chưa kịp nói hết câu thì mất thăng bằng, trượt thẳng về phía Flare. ``Gượm đã-!''

*RẦM!*

Cú va chạm làm nàng Dị giới chao đảo, còn Marine thì ``Á\~{}'' một tiếng đầy\dots\ đáng ngờ. Cô nàng hải tặc bắt đầu lộ rõ bản chất ``hỏny'' (dê xồm) của mình, nhất là khi mắt cô bắt đầu lấp lánh hình trái tim.

Nàng Thuyền trưởng lao tới với vẻ mặt sung sướng: ``Ufufu\~{} Matsuri-senpai à\~{}''

Nàng Lễ hội nhìn bản mặt của đàn em, hoảng sợ luống cuống quay đầu lại, vô tình vớ lấy sợi dây thừng gần đó.

*ỤP!*

Marine ôm chầm lấy đàn chị, đẩy ngã cô ấy xuống, kéo cả cái dây kia.

*TÁCH!*

Cả căn phòng bỗng im lặng trong một giây ngắn ngủi.

Mio, Aqua và Kanata đồng loạt ngước lên trần nhà. Một linh cảm chẳng lành ập đến. Và rồi\dots

*BỐP!* *COONG!*

Một chiếc chậu nhôm từ đâu rơi thẳng xuống đầu Kanata. Mio giật mình với tiếng động đó, cụp cả tai xuống, trong khi Aqua thì không biết nên hoảng hốt hay bật cười.

\begin{figure}[H]
  \centering
  \begin{subfigure}[t]{0.45\textwidth}
    \centering
    \includegraphics[width=8cm]{../../../../Image/Book?/7-11i1.png}
    \caption{Marine vô tình va chạm vào Flare, lao tới Matsuri\dots}
  \end{subfigure}
  \quad
  \begin{subfigure}[t]{0.45\textwidth}
    \centering
    \includegraphics[width=8cm]{../../../../Image/Book?/7-11i2.png}
    \caption{\dots khiến cô nàng hoảng hốt nắm lấy sợi dây thừng lủng lẳng\dots}
  \end{subfigure}
  \quad
  \begin{subfigure}[t]{0.45\textwidth}
    \centering
    \includegraphics[width=8cm]{../../../../Image/Book?/7-11i3.png}
    \caption{\dots làm chiếc chậu nhôm rơi thẳng vào đầu của Kanata.}
  \end{subfigure}
  \caption{Tóm tắt sự việc hỗn loạn của nhóm năm cô gái.}
\end{figure}

Nàng Sói đen thở dài ngán ngẩm, thở dài không phục: ``Được rồi\dots\ Tiếp theo!''

\ct

Tiếp tục là hàng loạt ý tưởng khác - mà đúng hơn là một chuỗi những trò quái đản không rõ có giúp họ trở thành idol tốt hơn hay không.

\il

Đầu tiên, họ thử nghiệm trò ``What's in the box?'' (Trong hộp có gì?), nhưng thay vì đặt đồ vật bí ẩn vào hộp, Marine và Matsuri lại\dots\ quyết định bịt mắt rồi sờ mặt Kanata.

\img{7-11j1}

Kanata, người bị chỉ định ở bên trong hộp, rùng mình khi thấy có hai bàn tay đang sờ soạng khắp mặt: ``Này! Cái đó là lỗ mũi của em mà!''

Flare chỉ đứng khoanh tay nhìn, không khỏi thắc mắc liệu trò này có dính dáng gì đến idol. Còn Mio thì chỉ biết xoa trán thở dài: ``Tiếp theo!''

\il

Sau đó, họ chơi trò ``Bốc bánh may mắn'' theo phong cách Towa-sama - một trò đầy rủi ro khi có một chiếc bánh chứa đầy mù tạt. Aqua, Matsuri và Mio may mắn an toàn. Nhưng Kanata\dots

``!\$)*\$\#(\$\#!\$\~\~\~!!!'' nàng Thiên sứ rên la trong đau đớn, với nước mắt giàn giụa, khuôn mặt đỏ bừng vì cay.

\img{7-11j2}

Nàng Sói đen dù thở phào nhẹ nhõm vì không bị dính bẫy mù tạt lần nữa\footnote{Xem lại {\flaregothic ÉPISODE 79}, phần {\abv Truyện phụ}, quyển số Sáu.}, không chút thương xót mà nói: ``Tiếp theo!''

\il

Kế đến, Marine và Aqua nảy ra một trò còn tàn nhẫn hơn: lấy vầng hào quang hình ngôi sao bốn cánh trên đầu Kanata để chơi trò ném và bắt.

\img{7-11j3}

Nàng Hầu gái cười toe toét, ném cái hào quang về đối phương như một cái đĩa: ``Nào, Marine-chan! Bắt lấy!''

Nàng Thuyền trưởng giơ tay chờ sẵn, chụp lấy cái ``của quý'' của cô nàng Thế hệ thứ Tư: ``Được, để đó cho em!''

Ở giữa, nàng Thiên sứ hốt hoảng, cố chụp lấy cái vầng hào quang của mình, gần như phát khóc: ``\textbf{Không được! Trả nó lại cho em!!}''

\il

Và cứ thế, các trò ngớ ngẩn tiếp tục nối tiếp nhau: từ nhảy dây bằng băng đô của Flare, đến việc Matsuri thử khiến Aqua hát khi đang uống nước, rồi Marine bày ra màn biểu diễn idol\dots\ nhưng chỉ toàn là tư thế tạo dáng kì quặc.

Từ ngoài cửa sổ, người ta có thể nghe thấy đủ thứ âm thanh hỗn loạn - tiếng la thất thanh của Kanata, tiếng thúc giục của những kẻ đang quá mức hào hứng, và cả những tiếng cười giòn tan đầy tinh quái, cùng tiếng thở dài ngao ngán của Mio.

\begin{dialogue}
  \speak{Kanata} *tuyệt vọng* \textbf{Dừng lại đi! Nghiêm túc đấy, mấy chị đang vượt quá giới hạn rồi!!}
  \speak{Marine} *cười đểu* Ngậm miệng đi! Cô em là ai?
  \speak{Kanata} *mếu máo* Là idol!
  \speak{Marine} *vỗ tay* Và chúng ta đang ở đâu?
  \speak{Kanata} *run rẩy* \textit{hololive}\dots
  \speak{Marine} *chỉ tay lên trời* Nghĩa là chúng ta chả có gì phải lo, đúng không!? \textbf{Đúng không!?}
  \speak{Flare} *giật mình* Marine, dừng lại đi!
  \speak{Matsuri} *chống nạnh* Chúng ta đang tập luyện để trở thành những idol tốt hơn đó!
  \speak{Aqua} *cuống cuồng* Mọi người, lùi xuống nào!
  \speak{Mio} *hoảng loạn* \textbf{Ai đó nắm lấy cái dây thừng đi chứ!}
\end{dialogue}

Và họ cứ thế tranh cãi, chất vấn nhau, rồi tiếp tục thực hiện những trò kì quặc đó đến tận chiều muộn. Nhưng cuối cùng, chẳng có cái nào trong số đó thực sự giúp họ cải thiện kĩ năng thần tượng như mong muốn cả.

\ct

Hoàng hôn đã buông xuống.

Tất cả đều đã kiệt sức. Flare ngả đầu ra sau trên ghế sô-pha, mắt lơ đãng nhìn lên trần nhà. Marine, vì một lý do nào đó, bị trói tay và đang ngồi bệt dưới sàn, thở dốc. Kanata thì nằm bẹp không nhúc nhích, như thể linh hồn đã rời khỏi thể xác. Matsuri khuỵu xuống, đầu gục về phía trước, lẩm bẩm với vẻ thất vọng: ``Không được rồi\dots\ Chúng ta đang luyện để trở thành diễn viên hài chứ có phải idol đâu!'

\img{7-11k}

Nàng Thuyền trưởng rên rỉ, tỏ vẻ thất vọng: ``Nếu cứ thế này, chúng ta sẽ chết mất thôi!''

Nói xong, Senchou cũng buông mình xuống sàn, mắt nhắm nghiền. Kanata chỉ có thể phát ra một tiếng thở dài chua chát, chẳng buồn phản ứng nữa.

Trên ghế sô-pha, Aqua và Mio cũng chẳng khá hơn. Cô nàng Sói đen thậm chí còn gục hẳn đầu xuống, đôi tai cụp lại đầy chán chường.

Không ai nói với ai câu nào. Một bầu không khí im lặng bao trùm khắp văn phòng. Dường như tất cả mọi người đều đã chịu bỏ cuộc, chấp nhận việc họ có thể sẽ bị sa thải vào một ngày không xa.

Ngoại trừ một người.

Một cô nàng Hầu gái với biệt danh ``Baqua'', người đang ngẫm lại những gì đã xảy ra trong ngày, cũng là người hăng hái nhất cho buổi huấn luyện ``có mà không'' này.

Aqua ngồi đó, đôi mắt tím long lanh phản chiếu ánh đèn mờ nhạt trong phòng. Cô nhìn xuống lòng bàn tay mình, rồi bất giác siết nhẹ. Một cảm giác mơ hồ len lỏi trong tâm trí, một điều gì đó mà cô đã luôn nghĩ đến nhưng chưa từng thực sự diễn đạt thành lời.

Cuối cùng, để phá tan bầu không khí sầu não, cô lên tiếng: ``Mà\dots\ cái gì mới làm nên một idol nhỉ?''

Câu hỏi bất chợt của cô làm cả phòng lặng đi. Mio khẽ nhấc đầu lên, ánh mắt đăm chiêu. Kanata liếc nhìn Marine, còn Flare thì ngả người về phía trước, hai tay đan vào nhau.

Ba người đang ``ôm'' sàn lần lượt ngước nhìn cô, rồi đáp lại theo bản năng, nhưng trong giọng họ vẫn có chút mơ hồ, như thể họ cũng chưa thực sự chắc chắn về câu trả lời của mình.

\begin{dialogue}
  \speak{Matsuri} Ờm\dots\ Khi ta tỏa sáng trên sân khấu sao?
  \speak{Kanata} Khi ta đẹp và giỏi như thiên thần?
  \speak{Marine} *cười yếu ớt* Khi ta ``cứ cố gắng làm'' thôi?
\end{dialogue}

Aqua im lặng, để những lời đó trôi qua như làn gió thoảng. Chúng đều có lý, nhưng sâu trong lòng, cô biết mình đang tìm kiếm một điều gì đó khác.

Cô cúi đầu, hai bàn tay đặt lên đầu gối, siết nhẹ rồi thả lỏng. Đôi vai nhỏ bé run rẩy một chút trước khi cô hít vào một hơi thật sâu.

Cuối cùng, cô hít một hơi thật sâu và cất tiếng, giọng có chút nghẹn ngào: ``Với chị thì\dots''

\img{7-11l}

Cô ngừng lại một lúc, như thể đang cân nhắc từng câu chữ. Flare hơi nhướng mày, còn Kanata quay hẳn đầu sang nhìn cô.

Đến khi đã bình tâm lại, Aqua tiếp tục: ``Chị nghĩ rằng idol là một người tỏa sáng và mang đến những nụ cười cho khán giả, trong khi họ cống hiến hết mình cho những gì họ yêu quý.''

Matsuri chớp mắt, dường như nhận ra điều gì đó. Cô đã từng nghe những lời như vậy, nhưng khi chúng được nói ra từ Aqua - một người luôn vụng về nhưng chưa bao giờ từ bỏ - chúng bỗng trở nên nặng hơn, sâu hơn.

``Mình không quan tâm người ta nói rằng những gì mình làm chẳng giống một idol tí nào cả.'' cô tiếp tục, siết nhẹ vạt váy.

Marine hơi mở miệng, nhưng lại chẳng nói gì. Cô biết Aqua đã từng đấu tranh rất nhiều với danh xưng ``idol'', đặc biệt là khi nhiều người ngoài kia đặt ra những tiêu chuẩn khắt khe đến mức phi lý.

Với mình, cách tỏa sáng của mình chỉ đơn giản là chơi game và nói chuyện với mọi người trên stream mà thôi,'' cô mỉm cười, nhưng nụ cười ấy lại có chứa một chút sự suy tư. Những giọt lệ rơi thẳng xuống đôi tay nhỏ nhắn của cô.

Mio ngẩng đầu lên, ánh mắt thoáng dao động. Dù có đôi lúc trêu chọc Aqua vì sự vụng về của cô, nhưng giờ đây, nàng Sói đen hiểu rằng: đó không phải là sự thiếu cố gắng, mà là cách riêng để Aqua kết nối với thế giới này.

Cô nắm hai bàn tay lại, tiếp tục với dòng tự sự của mình: ``Mình muốn làm những điều mà mình thích, cùng với tất cả mọi người trong fan và các thành viên \hll\dots''

Cô dừng lại, đôi mắt ánh lên sự chân thành tuyệt đối. Cô khẽ khàng, đôi mắt vẫn còn long lanh, nhưng giờ nó có toát ra một vẻ kiên định: ``Và đó mới là kiểu idol mà mình muốn trở thành.''

Sự im lặng bao trùm căn phòng.

Không ai nói gì, nhưng ánh mắt của họ dần thay đổi. Không còn sự chán chường hay mệt mỏi, mà thay vào đó là sự thấu hiểu và đồng cảm.

Lời bộc bạch của cô vang lên trong căn phòng tĩnh lặng, như một ngọn đèn nhỏ thắp sáng bóng tối. Từng câu chữ đơn giản nhưng đầy chân thành, khiến tất cả mọi người đều lắng nghe một cách nghiêm túc. Và rồi, họ dần nhận ra: những điều Aqua vừa nói, đó là sự thật.

Flare khẽ tựa cằm lên tay, khóe môi hơi nhếch lên. Kanata chớp chớp mắt, như thể đang cố kìm nén điều gì đó trong lòng. Matsuri ngả người ra sau, mắt nhìn lên trần nhà, miệng khẽ cười như thể cô vừa nhận ra một điều gì đó quan trọng.

Marine nhắm mắt, rồi thở ra một hơi dài Cô nhoẻn miệng cười: ``Thật là\dots\ Akutan lại khiến mọi người rơi vào bầu không khí nghiêm túc này\dots''

Mio bật cười nhẹ, nhưng trong mắt cô vẫn còn đọng lại chút gì đó sâu sắc. ``Có lẽ Aqua nói đúng.''

Không ai phản bác. Vì tất cả bọn họ đều hiểu - idol không phải chỉ là đứng trên sân khấu, không phải chỉ là những hào quang chói lòa.

Mỗi người đều có cách tỏa sáng của riêng mình.

Và cách mà Aqua chọn\dots\ cũng chẳng hề sai.

Mio nhìn Aqua thật lâu trước khi mỉm cười, như thể cô vừa tìm thấy một câu trả lời cho chính mình. Cô nắm lấy tay người đàn chị, nở một nự cười dịu dàng: ``Mọi người ai cũng có cái lý về idol cho riêng mình. Họ biết họ muốn trở thành một idol như thế nào sao cho vừa lòng với khán giả nhất mà.''

Một câu nói không cần phải nói to, nhưng cũng vừa đủ sức nặng để bất cứ ai nghe được cũng phải thực sự suy nghĩ rất lâu. Vì đó là sự thật. Idol không phải một khuôn mẫu cố định mà ai cũng phải đi theo - nó là hành trình mà mỗi người tự vẽ ra cho mình.

Matsuri ngả người về phía trước, ánh mắt sáng lên. ``Đúng rồi đó! Việc phán xét và đuổi việc một người nào đó mà chỉ dựa theo những thứ đó thôi quả là điều không công bằng!''

Cô siết chặt tay lại, rõ ràng vẫn còn bức xúc về tình cảnh mà họ đang gặp phải. Marine cũng gật gù đồng tình, đặt một tay lên hông: ``Vậy ngày mai chúng ta hãy cùng nói chuyện với cấp trên của ta nha!''

Flare mỉm cười, ánh mắt mang theo chút tự tin đầy mạnh mẽ. ``Tôi đồng ý! Nghe có vẻ đây sẽ là kế hoạch cho chúng ta nhỉ?''

Kanata siết chặt nắm tay. ``Nếu ta không thể nói chuyện được với Kuon-senpai, vậy ta có thể trao đổi với A-chan một lần nữa, thậm chí là YAGOO nữa!''

``Để cho mọi việc ra ngô ra khoai luôn!'' tất cả mọi người cùng nhau giơ nắm đấm lên trời, ánh mắt hiện rõ một sự quyết tâm cao độ.

Không còn bầu không khí chán chường hay thất vọng như lúc trước nữa. Sự bùng cháy quyết tìm ra sự thật dần hiện rõ trên gương mặt của từng người.

\img{7-11m}

Aqua nhìn họ với sự ngạc nhiên, đôi mắt tím tròn xoe. Chưa bao giờ cô nghĩ rằng chỉ với những suy nghĩ chân thành của mình lại có thể tạo ra động lực lớn đến vậy. Những người chị em bên cạnh cô, họ thực sự tin tưởng vào điều này.

Cô khẽ cúi đầu, môi mím lại như thể đang suy nghĩ gì đó. Nhưng khi ngẩng lên lần nữa, trên gương mặt nhỏ nhắn ấy đã nở một nụ cười nhẹ. Một nụ cười không cần lời nói, nhưng ai cũng hiểu rằng cô đã đồng ý.

Mio nhìn quanh, thấy ai cũng đã sẵn sàng. Cô mỉm cười, ngẩng đầu lên đầy tự tin. ``Vậy là mọi người theo ý kiến rồi chứ gì?''

Tất cả cùng gật đầu đồng tình với cô.

``Vậy chúng ta hãy đòi lại công bằng cho chính chúng ta nào! \textbf{Cố lên!}''

Kanata, Flare, Marine (giơ hai tay) và Matsuri cùng Aqua cũng nhiệt liệt hưởng ứng theo: ``\textbf{Cố lên!}''

Họ không còn chỉ là những idol đang bị đẩy vào thế bị động nữa. Họ đã sẵn sàng để đứng lên vì chính con đường mà mình đã chọn.

\begin{center}
  \uline{Sáng hôm sau.}
\end{center}

Trời vẫn còn hơi lạnh, ánh sáng buổi sáng nhàn nhạt chiếu qua những ô cửa kính của công ty. Những con đường đã tấp nập người qua lại chuẩn bị cho một ngày làm việc mới, nhưng nhóm sáu người lại không có tâm trạng để hưởng thụ bầu không khí đó.

Họ bước vào công ty với sự quyết tâm hiện rõ trên từng gương mặt. Mio dẫn đầu, theo sau là Flare, Kanata, Matsuri, và lần này cả Aqua lẫn Marine cũng đi cùng.

Nàng Lễ hội, như thường lệ, là người chủ động nhất. Cô đẩy mạnh cánh cửa văn phòng, giọng to như sấm: ``\textbf{Quản lý có ở đây không ạ!?}''

\img{7-11n}

Đáp lại tiếng gọi ấy chỉ là một sự im lặng đến lạ thường.

Không ai trong nhóm có thể thấy bất kỳ nhân viên nào trong phòng cả, không một tiếng gõ phím, không một lời trò chuyện nào vang lên.

Nàng Thiên sứ nheo mắt nhìn quanh, lẩm bẩm: ``Em không thấy ai trong phòng cả.''

Nàng Yêu tinh vàng khoanh tay lại, nghi hoặc: ``Họ đang lơ đễnh rồi sao? Có phải là như Kuon-senpai đâu\dots''

Nàng Hầu gái, vốn ít khi chủ động, đột nhiên chỉ tay về phía bàn làm việc gần cửa sổ: ``Không đâu. A-chan đang ở kia kìa.''

Mọi người lập tức dõi theo ánh nhìn của cô. Ở góc văn phòng, một cô gái đeo kính đang chăm chú vào màn hình, tay gõ bàn phím liên tục. Cô ấy có vẻ đang chuẩn bị một bài thuyết trình trên web.

\img{7-11o}

Họ cùng thở phào nhẹ nhõm. Ít nhất, có vẻ như cô ấy đã quay trở lại làm việc như bình thường.

Vừa nghe thấy tiếng động, nàng Đeo kính liền quay đầu lại, nở một nụ cười nhẹ: ``Ồ, chào buổi sáng mọi người.''

Nhưng trước khi cô kịp nói gì thêm, Aqua đã chạy thẳng đến chỗ cô, khuôn mặt đầy nghiêm túc: ``A-chan à!''

A-chan giật mình trước sự nhiệt tình bất ngờ này: ``Sao?''

Nàng Hầu gái hít một hơi sâu rồi nói, giọng điệu rõ ràng và đầy quyết tâm: ``Có thể bọn em vẫn còn xa lắm mới có thể trở thành một idol hoàn hảo được\dots''

Mio bước lên, chậm rãi tiếp lời: ``\dots nhưng chính bọn em mới là người quyết định bọn em sẽ trở thành một idol như thế nào mà.''

Flare gật đầu, đôi mắt sắc bén thể hiện rõ sự tự tin: ``Cùng với các fan hâm mộ của chúng em, chúng em sẽ tiến gần hơn đến mục tiêu.''

Matsuri siết chặt tay: ``Và chúng em nghĩ sẽ thật là bất công nếu một người nào đó đánh giá và sa thải chúng em chỉ vì mấy thứ đó ạ.''

Marine bước lên, hai tay ôm ngực, giọng đầy bi thương như thể đang diễn một vở kịch buồn: ``Bọn em muốn chị cân nhắc lại mấy cái đánh giá và sa thải này ạ. Nếu không thì\dots''

Cô dừng lại một giây, rồi đột ngột đổi giọng đầy đau khổ, quỳ xuống một cách thảm thương: ``Thuyền của em\dots\ Không, giấc mơ của em sẽ bị chìm nghỉm ạ!''

Kanata theo phản xạ cũng bật ra một câu đầy lo lắng: ``Nợ của em phải trả nữa!''

A-chan im lặng, trấn tĩnh lại sau những câu nói dồn dập từ sáu cô giá. Cô đẩy gọng kính lên, ánh mắt dường như chưa kịp hiểu hết mọi chuyện.

Cuối cùng, cô thở ra nhẹ nhàng rồi đáp: ``Thì, công ty chúng ta không có định sa thải bất cứ ai đâu, đấy là điều đầu tiên chị muốn nói. Kể cả Hikawa-san cũng thế.''

Mọi người đứng hình.

Nàng Đeo kính tiếp tục, giọng điềm tĩnh nhưng không kém phần dịu dàng: ``Và bọn chị cũng nghĩ các em phát triển thành những idol riêng của bản thân mình cũng là một điều đáng quý mà.''

Aqua nhíu mày, chớp mắt vài lần trước khi cất giọng, như thể muốn chắc chắn rằng mình không nghe lầm: ``Đúng chứ!?''

Nhưng trước khi bất cứ ai nói tiếp điều gì, tất cả mọi người đã đồng loạt thốt lên trong sự kinh ngạc: ``Hả!?''

\img{7-11p}

A-chan nhìn họ, ngơ ngác chẳng kém: ``Ủa?''

Sự im lặng lại bao trùm căn phòng trong vài giây. Đây không phải là kết quả họ mong đợi. Tất cả sự chuẩn bị tinh thần cho một cuộc tranh luận nảy lửa\dots hóa ra lại hoàn toàn vô dụng.

Kanata cầm tờ báo cáo lên, vẫn chưa tin vào tai mình: ``Nhưng chính chỗ này ghi rằng chúng ta đang có kế hoạch cắt giảm mà!''

Cả nhóm lập tức quay sang nhìn tờ giấy trên tay Thiên sứ, như thể nó đang cầm giữ vận mệnh của họ.

Nàng Đei kính nhẹ nhàng cầm lấy tờ giấy, điều chỉnh lại kính rồi chăm chú đọc qua từng dòng chữ. Cô im lặng một lúc, ánh mắt lướt qua những đoạn văn bị tẩy xóa, rồi chợt ``À!'' lên một tiếng như vừa nhớ ra điều gì đó.

Cô vỗ nhẹ vào tờ giấy rồi giải thích: ``Mấy đứa có biết cái này là gì không? Cái này là bản tuyên ngôn ăn kiêng của một trong những nhân viên ở đây đấy.''

Không khí trong phòng bỗng chốc\dots\ lặng như tờ. Một cơn gió lạnh vô hình thổi qua, và ngay sau đó\dots

``\textbf{Hảaaaaa!???}'' cả sáu người đều bật ngửa vì cú sốc. Một số ôm đầu, một số đập bàn, một số chỉ đơn giản là há hốc mồm.

Matsuri, mắt giật giật, giơ tay lên trước: ``Thế\dots\ con số 100 là gì?''

A-chan nhún vai, đáp lại tỉnh bơ: ``Có lẽ đó là số cân mà người đó đã vượt quá.''

Mio, vẫn còn chưa hoàn hồn, hỏi tiếp: ``Còn hạng D thì sao?''

``Chị tin là từ kết quả khám sức khỏe của người đó.''

Flare, cố giữ vẻ bình tĩnh nhưng cô đã đổ mồ hôi hột: ``Và\dots\ loại bỏ cái gì?''

``Mỡ thừa.''

Một tiếng ``phụt!'' vang lên khi Marine suýt nữa phun nước ra. Nhưng cô vẫn cố gắng hỏi tiếp với vẻ mặt như không tin vào thực tế trước mắt: ``Vậy tại sao lại có đoạn bị tẩy xóa với bị che vậy?''

Nàng Đeo kính bật cười, một nụ cười đầy ẩn ý: ``Chị nghĩ chắc người đó cảm thấy ngượng sau khi viết nó ra đấy.''

\dots

Im lặng.

Một khoảng lặng dài, rất dài.

Chỉ có tiếng tích tắc của đồng hồ và hơi thở phập phồng của sáu con người vừa nhận ra rằng -

Họ đã bị ăn một cú lừa rất mạnh.

Không có ai bị sa thải cả. Không có một cuộc cắt giảm nhân sự nào hết.

Tất cả những lo lắng, căng thẳng, kịch bản đấu tranh\dots\ đều xuất phát từ một bản kế hoạch giảm cân vô tình bị hiểu lầm.

Và rồi\dots

``\textbf{Đúng là lừa gạt mà!}''

\img{7-11q}

Cả sáu người đồng thanh hét lên, mỗi người một cảm xúc khác nhau: người thì giận dữ đập bàn, người thì cười trừ vì thấy bản thân quá ngớ ngẩn, người thì chỉ biết thở dài vì đã mất một buổi sáng - hoặc thậm chí là mất nguyên ngày hôm qua - quý giá vào việc này.

A-chan chỉ đáp lại bằng một nụ cười chát, vừa có chút áy náy, vừa có chút thích thú trước phản ứng của họ.

Cô chỉnh kính một lần nữa rồi lẩm bẩm: ``Nhưng mà\dots\ tờ giấy này sao lại lạc sang văn phòng của chúng ta nhỉ?''

Nhưng ngay sau đó, cô chợt nhớ ra một chuyện quan trọng hơn: ``À, tiện thể các em đang ở đây, bọn chị đang có kế hoạch cho giai đoạn tiếp theo của dự án idol chúng ta này.''

Không khí trong phòng lập tức thay đổi. Cả sáu người đều tập trung lắng nghe, ánh mắt đầy háo hức.

A-chan tiếp tục: ``Ý tưởng lần này là lập các nhóm biểu diễn khác nhau và tổ chức một buổi concert. Chị nghĩ chắc Hikawa-san cũng đã nói qua với các em rồi.''

Vừa nghe đến đó, tất cả họ đều rạo rực. Một số gật gù, một số thì lập tức nhớ lại một chuyện quan trọng\dots

Nàng Sói đen suy nghĩ một chút rồi lên tiếng: ``Nhắc mới nhớ, Kuon-senpai cũng đã từng nhắc đến nó rồi\dots''

Nàng Lễ hội búng tay cái, chợt nhớ ra điều quan trọng: ``À đúng rồi! Lúc mà anh ấy soạn báo giá ấy!''

\begin{center}
  \uline{Hồi tưởng\\Một tháng trước.}
\end{center}

Trong phòng làm việc của văn phòng, Kuon đang ngồi trước máy tính, vẻ mặt đầy tập trung. Trên màn hình là một loạt bảng tính chi phí, và bên cạnh cậu là một máy tính cầm tay cùng vài tờ tài liệu.

Cậu vừa gõ bàn phím, vừa lẩm bẩm tính toán: ``\vie{Phí xây dựng sân khấu từ Khoa Kiến trúc công trình\dots\ 100 nghìn yên\dots\ Rồi phí trang trí gồm vẽ họa tiết, sơn\dots\ tất cả hết khoảng 40 nghìn yên\dots}''

Rồi cậu bất giác rên lên một tiếng: ``\vie{Mẹ vãi chưởng! Mất từng ấy tiền cơ đấy!}''

Cậu thở dài, tay chống cằm: ``\vie{Huy động cả trường làm mà đã mất từng này, chả biết thuê ngoài còn mất bao nhiêu nữa\dots\ Mà cô-vít, cô-veo kiểu này thì phải tiết kiệm thôi\dots}''

Đúng lúc ấy, từ phía sau, một giọng nói bất thình lình vang lên: ``Anh đang tính toán gì đấy?''

``\textbf{Ối giồi ôi!!}'' Kuon giật mình suýt ngã khỏi ghế, vội quay lại nhìn.

Cậu thấy Matsuri đứng đó, tay chống hông, mặt đầy tò mò.

``C-Có việc gì à?'' cậu hỏi, cố gắng tự trấn an bản thân.

Nàng Lễ hội cười lớn: ``Là em đây mà, Matsuri đây! Em thấy anh tính toán từ nãy tới giờ trông chăm chú lắm. Cái gì vậy?''

Cậu Tổng nhấc vài tờ báo giá lên, huơ huơ trước mặt: ``À, báo giá lắp đặt sân khấu với cơ sở vật chất cho buổi live thôi mà.''

Mio vừa lúc đi ngang qua, nghe thấy liền dừng lại: ``Buổi live vừa rồi của chúng ta à?''

Cậu lắc đầu: ``Không. Buổi live sắp tới. Ơn giời là làm online nên chi phí rẻ hơn hẳn.''

Nghe vậy, Matsuri trợn mắt: ``Thế á!? \hll\ chúng ta sắp có concert sao!?''

``Đã thế lại còn online nữa!?'' Mio cũng ngạc nhiên không kém.

Cậu Tổng cười nhẹ: ``Ờ. Tin chuẩn em nha. Ngoài ra còn nhiều thứ mới nữa. Chính YAGOO bảo anh soạn báo giá này đấy chứ.''

Cả hai cô gái càng tỏ vẻ tò mò hơn: ``Cụ thể là thế nào vậy?''

Cậu chỉ đáp lại bằng một cái lắc đầu: ``Tuyệt mật. Chờ đến hôm đó hẵng hay.''

Nàng Lễ hội nghe đến đây bĩu môi: ``Chán thế\~{}\dots''

Dù cho cậu Tổng không tiết lộ quá nhiều, nhưng chỉ bấy nhiêu thôi cũng khiến cả hai người họ cảm thấy vô cùng háo hức.

\begin{center}
  \uline{Kết thúc hồi tưởng.}
\end{center}

Sau khi Mio và Matsuri hồi tưởng lại khoảnh khắc ấy, A-chan gật đầu xác nhận: ``Chính chị với Hikawa-san đã khai triển buổi concert này, từ lập kế hoạch, đến họp bàn, đến xây dựng sân khấu đấy.''

Rồi, cô chợt cười, hơi tiếc nuối: ``Tiếc là cậu ấy không thể đến dự buổi live đó với các em được vì chuyện gia đình. Nhưng chị nghĩ ít nhất cậu ấy cũng đã đóng góp kha khá cho ngày sắp tới rồi.''

Aqua hào hứng giơ tay lên: ``Vậy là mọi thứ đã chuẩn bị xong rồi sao!?''

Cô nàng Đeo kính gật đầu: ``Có thể nói là gần hoàn thiện rồi. Giờ chỉ còn chia nhóm để biểu diễn thôi. Chúng ta sẽ chia làm nhiều nhóm thành nhiều đợt khác nhau, và ta cứ làm theo format của Fes.''

Cả nhóm càng nghe càng phấn khích. Một concert theo hình thức chia nhóm sẽ mang lại nhiều điều mới lạ, và tất cả họ đều muốn tham gia đóng góp.

\img{7-11r}

Matsuri hăng hái lên tiếng trước: ``Thế ạ? Nói thêm cho bọn em biết đi!''

Marine cười khoái chí: ``Wow, nghe có vẻ vui đấy.''

Aqua cũng giơ cả hai tay lên đồng tình: ``Em đồng ý!''

Kanata chợt nghĩ ra một ý tưởng mới: ``Em muốn, kiểu, làm những kịch hài nhỏ trong buổi concert đó ấy!''

Flare tán thành: ``Có một nhóm ba người sẽ giúp cho chúng em nghĩ ra được nhiều câu đùa hơn đó.''

Mio hơi nhăn mặt: ``Lại quay về hài rồi!?''

Mỗi người lại đưa ra một ý tưởng khác nhau, nhưng may mắn thay, lần này A-chan đã ghi chú lại tất cả một cách cẩn thận. Cô lắng nghe, phân tích và phản hồi rõ ràng về từng đề xuất - có cái được chấp nhận, có cái phải chỉnh sửa, và cũng có cái bị loại bỏ vì không phù hợp.

Nhưng điều quan trọng nhất là ai cũng cảm thấy vui vẻ. Bởi lẽ, tất cả đều đang hướng đến cùng một mục tiêu:

Mang đến một buổi concert thật tuyệt vời cho khán giả!

\pagebreak

\begin{center}{\ab BONUS SONG}\\(Đây đơn thuần chỉ là phần hát của các thành viên sẽ xuất hiện trong buổi live Bloom, đó. Bạn có thể xem video ở dưới, nếu không thích đọc phần tự sự của tôi.)
  \vid{7-11s}{song}\\(Đặt giả thuyết là tất cả mọi người đều cùng hát bài hát đó.)
\end{center}

Câu hát mở đầu với hai câu:

\begin{center}
  少しずつ近づいているよ　夢に見たトキメキ \\
  (Niềm háo hức mà chúng ta hằng mong chờ đang đến gần)

  広がってく世界中に \\
  (Tất cả, chúng được rải ra trên khắp thế gian này)
\end{center}

Cùng với đó, tất cả các thành viên hololive đều dần dần xuất hiện, bắt đầu di chuyển (dạng hình tĩnh) từ phải sang: Roboco, Mel, Aki, Matsuri. Câu hát tiếp tục:

\begin{center}
  皆を笑顔にさせるキラメキを \\
  (Chúng ta đây sẽ tỏa sáng lấp lánh, khiến cho mọi người đều mỉm cười)

  一人じゃないからライバルがいるから \\
  (Vì chúng ta không hề đơn độc, vì chúng ta đều có đối thủ của nhau)
\end{center}

Tiếp tục là Fubuki, Aqua, Shion và Mio.

\begin{center}
  競い合うことすら \\
  (Dù cho là một cuộc chiến khốc liệt đến cỡ nào)

  楽しくなっちゃうんだよ \\
  (Nó cũng sẽ trở thành một niềm vui)
\end{center}

Kế đến là Ayame, Subaru, Choco và Miko.

\begin{center}
  {\fotskip Dear My Friends} \\
  (Hỡi người bạn của tôi)

  未来へ走り出したら{\fotskip Fantastic} \\
  (Khi ta hướng tới tương lai, mọi thứ sẽ rất là tốt đẹp ở phía trước)

  夢を咲かせて{\fotskip Our Future} \\
  (Ta sẽ cùng biến những ước mong nở rộ của mình vào tương lai của chúng ta)

  一歩一歩を踏み出して　トップを目指そう \\
  (Từng bước, từng bước, ta từ từ hướng tới đỉnh cao)

  信じる力が心の{\fotskip Happiness} \\
  (Hãy cùng tin tưởng vào bản thân, để rồi sẽ đưa con tim mình trở nên hạnh phúc)
\end{center}

Tiếp theo là năm thành viên thế hệ thứ Ba: Pekora, Rushia, Marine, Flare và Noel.

\begin{center}
  全力で楽しもう \\
  (Ta cùng tận hưởng những gì mà ta có)

  それはいつかきっと \\
  (Bởi chắc chắn một ngày nào đó)
\end{center}

Cuối cùng là bốn thành viên mới nhất từ \hll\ thế hệ thứ năm: Botan, Polka, Lamy và Nene\footnote{Sẽ được giới thiệu ở quyển số Tám.}

\begin{center}
  大きな夢叶えるから \\
  (Chúng ta sẽ biến những ước mơ to lớn đó trở thành sự thật)

  {\fotskip Dreaming Days} \\
  (Hỡi những ngày mơ mộng của chúng ta\dots)
\end{center}

Đoạn phần ca hát đó kết thúc, màn hình quảng cáo cho Bloom, xuất hiện, cùng với đó là câu kết:

\begin{center}
  \textbf{Đây sẽ là nơi mà chúng ta nở rộ, ngay tại chính nơi này\dots}
\end{center}

%==========================================TRUYỆN PHỤ=========================================
\omk

\pov{Hikawa Kuon}

Mùng 7 Tết Âm lịch.

Đã mười ngày kể từ khi tôi trở về quê ăn Tết. Đây là cái Tết thứ hai kể từ lúc tôi nhận chức Tổng quản lý nhóm \hll.

Nếu tính theo lịch Dương, hôm nay cũng đã là một ngày sau buổi concert online \textit{Bloom,}. Đáng lẽ đó phải là một buổi diễn tràn ngập tiếng reo hò cổ vũ từ khán giả, nhưng tiếc thay, đại dịch đã khiến họ không thể biểu diễn tham dự trực tiếp.

Chính xác hơn, họ đã phải biểu diễn mà không có khán giả. Thiếu sự nhiệt tình của khán giả, cũng như không có sự tương tác với những người hâm mộ hẳn là một tổn thất lớn, nhưng với tình hình hiện tại thì chỉ còn cách phải tuân theo thôi.

Dù vậy, dự án này vẫn đánh dấu một bước tiến quan trọng của hololive - một dự án còn khá mới mẻ, nhưng đã nhận được sự đầu tư rất lớn cả về chất xám lẫn tài chính. May mắn thay, nhìn vào những phản hồi tích cực từ khán giả, có thể nói nó đã rất thành công.

Tất nhiên, thành công đó đến từ sự tập luyện chăm chỉ và cống hiến không ngừng nghỉ của các thành viên. Nhưng bên cạnh đó, tôi cũng không thể không cảm ơn những người bạn trong trường đại học của mình. Chính họ đã giúp sức rất nhiều để hoàn thiện phần hậu kỳ cho buổi biểu diễn này, không chỉ nâng cao chất lượng mà còn giúp tôi tiết kiệm được một khoản đáng kể.

Thế nhưng, trớ trêu thay, tôi lại chẳng thể theo dõi buổi hòa nhạc ấy một cách trọn vẹn. Tôi không có mặt ở hậu trường dù thật lòng mà nói, tôi cũng không rõ liệu mình có thể tham dự không. Ngay cả việc xem trực tiếp qua điện thoại cũng bất khả thi, bởi nhà tôi\dots\ đang thực hiện chính sách tiết kiệm chi phí một cách khá triệt để.

Những gì tôi biết về buổi diễn chỉ là qua các tin tức đọc được trên mạng. Tôi có thể hình dung được sự thất vọng của mọi người khi biết tôi đã không theo dõi \textit{Bloom,} như họ mong đợi. Dù vậy, khi nghe A-san và đặc biệt là Tanigo-san (YAGOO) báo tin mừng, tôi cũng đủ hình dung ra nó đã thành công đến mức nào. Tôi gửi lời chúc mừng, cảm ơn mọi người vì sự cống hiến và cũng không quên xin lỗi vì đã bỏ lỡ giây phút quan trọng ấy. Nhưng dù thế nào đi nữa, cuối cùng thì bao nhiêu nỗ lực của chúng tôi cũng đã được đền đáp. Với tôi, chỉ cần vậy là đủ để cảm thấy mãn nguyện rồi.

Họ vui mừng là thế, còn tôi thì sao? Bình thường thôi. Cũng giống như mọi cái Tết mà tôi đã trải qua suốt hơn hai mươi năm nay.

Gọi là nghỉ Tết, nhưng thực chất cũng chỉ là danh nghĩa. Từ lúc tôi về đến hết mùng 5, tôi tạm gác công việc sang một bên. Nhưng ngay sau đó, tôi lại lao vào kế hoạch tiếp theo cho \hll\ - chủ yếu là trao đổi trực tuyến với A-san và Tanigo-san. Nếu có ai đó nghĩ tôi là người chăm chỉ\dots\ cũng không hẳn. Vì thật ra, tôi có lý do riêng để làm vậy.

Nói ra thì có vẻ như đang kể khổ, nhưng với tôi, ``ngày nghỉ'' là một khái niệm xa xỉ. Không phải vì công ty bắt ép, mà là do tôi tự muốn như vậy. Đến mức mà ngay cả Tanigo-san - một CEO dày dạn kinh nghiệm - cũng phải lắc đầu vì tôi ``tham việc quá'' dù ông ấy hiểu rằng trong văn hóa làm việc Nhật Bản, chuyện tăng ca chẳng có gì lạ lẫm.

Trong suốt một tuần, tôi chỉ nghỉ giỏi lắm là Chủ Nhật, còn các ngày lễ lớn, như Tết Dương, thì có thể miễn cưỡng dành ra một ngày. Những dịp mà tôi cho là ``vặt vãnh'' như Tuần lễ Vàng, Ngày Hiến pháp hay Ngày Thiếu nhi, tôi vẫn xung phong làm việc, chẳng bận tâm A-san hay Tanigo-san nghĩ gì về mình.

Làm việc là để kiếm sống, chẳng phải ai cũng vậy hay sao?

Những ngày mà lẽ ra tôi có thể nghỉ ngơi, cộng với gần một tháng nghỉ phép, cuối cùng cũng đủ để tôi an tâm về quê ăn Tết cùng gia đình. Dĩ nhiên, tôi vẫn được trả lương đầy đủ, thế nên ít ra cũng không cần lo lắng về tài chính. Và giờ thì, như các bạn thấy đấy, tôi đang ở nhà đón Tết.

Nhưng nói thật, Tết năm nay\dots\ hơi chán. Nếu so với năm ngoái - khi gia đình tôi đón hơn hai mươi người đến chung vui, tạo nên một không khí náo nhiệt không khác gì một hội chợ - thì năm nay rõ ràng ảm đạm hơn hẳn. Nhưng lạ một điều là tôi lại thích sự ``ảm đạm'' này. Nó yên tĩnh, không quá ồn ào, không phải tất bật tiếp khách hay lo toan đủ thứ. Chỉ có điều\dots\ tôi vẫn cảm thấy thiếu thiếu một chút gì đó.

Không phải vì tôi chán sự lặp lại của ngày Tết, mà chỉ là\dots\ không có những người đó, không khí có vẻ trống vắng hơn thường lệ. Nhưng mà thôi, tôi còn than thở gì chứ? Chiều tối mai hoặc chậm nhất là sáng ngày kia, tôi đã phải lên đường trở lại thành phố Điện tử để tiếp tục công việc rồi. Vẫn còn một nhiệm vụ quan trọng nữa: bàn giao người cuối cùng của \textit{NePoLaBo} - Momo\dots\ ờm, gì đó, Nene.

Nếu các bạn thắc mắc tại sao năm nay tôi không mời nhóm \hll\ đến ăn Tết như năm ngoái, thì có ba lý do:

\begin{itemize}
  \item Thứ nhất, gia đình tôi đang thực hiện chính sách tiết kiệm chi phí một cách khá nghiêm ngặt.
  \item Thứ hai, các idol-kiêm-hề cũng đang phải tập luyện gấp rút cho buổi biểu diễn, không có nhiều thời gian để nghỉ ngơi hay đi xa.
  \item Và quan trọng nhất, nhà tôi năm nay có đám giỗ của ông nội - người mất vào Rằm tháng Giêng năm ngoái. Theo phong tục ở đây, trong năm đầu tiên sau khi mất, gia đình sẽ không mời người ngoài đến nhà để tránh mang xui xẻo đến cho họ.
\end{itemize}

Còn một lý do nữa, tuy không phải là chính, nhưng chắc chắn có sức nặng. Đó là Mano Aloe-chan đang sống chung với chúng tôi - các bạn đã đọc ở quyển D thì chắc các bạn sẽ hiểu. Nếu các thành viên \textit{hololive} mà kéo đến thì ``kế hoạch của chúng ta coi như hỏng bét!'', hoặc ít nhất đó sẽ là lời của YAGOO. Thực ra tôi cũng không rõ cái kế hoạch đó cụ thể là gì, nhưng nếu nó đổ bể thì chắc chắn không phải là chuyện hay ho gì cả.

Nhắc đến nàng succubus ấy, mấy ngày vừa qua, em ấy vẫn đang cùng chúng tôi đón Tết. Mọi thứ cũng giống như những năm trước thôi, chẳng có gì quá đặc biệt để kể.

Ban đầu, tôi dự định sẽ nghỉ đến hết Rằm để dự đám giỗ gia đình, nhưng có lẽ tôi sẽ quay trở lại thành phố sớm hơn để tiếp tục công việc. Một guồng quay mới lại bắt đầu, và dự cảm của tôi cho thấy, năm nay tôi sẽ còn bận rộn hơn cả hai năm trước cộng lại.

Nhưng đó là chuyện của tương lai, còn bây giờ thì sao?

\il

Quay trở về thực tại.

Tôi cũng đang viết báo cáo cho công ty, giống như mọi hôm. Ngoại trừ\dots\ không hoàn toàn như vậy, vì lần này tôi đang phải cố hoàn thành nó trong khi xung quanh chẳng có lấy một chút không khí làm việc. Mọi người trong nhà vẫn còn tận hưởng những ngày Tết cuối cùng, còn tôi thì dán mắt vào màn hình, cố gắng hoàn thành nốt công việc trước khi trở lại thành phố Điện tử.

Trong lúc tôi đang cặm cụi gõ phím, bỗng màn hình máy tính báo có cuộc gọi video đến. Nhìn vào màn hình, tôi nhận ra ngay người gọi: giám đốc Tanigo-san.

``\vie{Chắc là gọi hỏi về công việc đây mà,}''  tôi lẩm bẩm, khẽ thở dài rồi bấm nhận cuộc gọi.

``À, cậu Hikawa-san đấy hả?'' Đúng, là ông ấy rồi.

Tôi vẫy tay chào lại: ``Vâng, chào chú ạ.''

``Mọi việc thế nào rồi?'' Đúng, tôi nghĩ là ông ấy đang hỏi tôi về công việc đây.

``À, cái bản báo cáo cháu soạn sắp xong rồi, tầm tối thì cháu-''

``Không, không, tôi không hỏi công việc. Tôi hỏi việc nhà cậu ấy.''

Tôi hơi bất ngờ. Bình thường, mỗi khi chú Tanigo-san gọi thì hầu hết là để bàn về công việc, chứ hiếm khi lại hỏi han chuyện cá nhân như thế này.

Dù tôi cảm thấy hơi nghi ngờ, nhưng tôi vẫn nhún vai trả lời: ``Nhà cháu ạ? Nhà cháu bình thường. Mọi người đều khỏe.''

``Ăn Tết vui vẻ chứ?''

``\dots Vâng\dots? Chú cũng có thể nói như vậy.'' Tôi trả lời một cách chung chung. Đúng là Tết năm nay có hơi tĩnh lặng hơn so với năm ngoái, nhưng nếu bảo tôi đánh giá nó có vui hay không thì cũng khó nói.

``Cậu nói thế là tôi thấy yên tâm rồi.''

``Thế ạ? Vậy để cháu cúp máy-'' tôi di chuột về nút kết thúc cuộc gọi, nhưng ông ấy chợt thốt lên: ``Ấy ấy, từ từ, cái cậu này! Đã nói xong đâu!''

``\dots Dạ?'' tôi nghiêng đầu khó hiểu.

``Tôi gọi cho cậu đâu chỉ mỗi hỏi thăm sức khỏe đâu. Họ cũng muốn nói chuyện với cậu đấy.''

``Họ? Ý chú là\dots''

Nghe có gì đó không ổn lắm. Có lẽ ông ấy đang không nói về những người đã cố gắng gọi cho tôi lúc tôi đang nghỉ Tết đấy chứ?

``Cậu đoán đúng rồi đấy. Đã lâu ngày rồi cậu chưa về, các cô gái giờ cũng nhớ cậu lắm.''

Tôi cau mày. ``Chú nói thế nào chứ trông họ vẫn vui mà. Ở trên live show họ diễn hôm qua ấy, cháu biết thừa.''

``Bên ngoài khác với bên trong đấy, cậu bé à. Thôi, để tôi dẫn họ nói chuyện với cậu.''

Nói rồi, ông giám đốc rời khỏi chỗ ngồi, màn hình rung nhẹ khi ông ấy di chuyển. Tôi không biết mình có nên vui hay lo lắng nữa\dots

Tôi nghe thấy tiếng ông ấy gọi mấy người vào. Trong lúc đó, tôi tranh thủ rời khỏi chỗ ngồi để lấy cốc nước, vừa để tránh bị gọi ngay lập tức, vừa để có thời gian chuẩn bị tinh thần.

Mà thực ra, cũng có gì đâu mà phải chuẩn bị?

Tôi không thể gọi đó là cảm giác khó chịu, nhưng thú thực là tôi vẫn chưa quen với những cuộc trò chuyện thế này. Một phần vì tính tôi không thích chen chúc trong những nơi đông người - live concert dù là online cũng không ngoại lệ. Đó là lý do tôi quyết định nghỉ Tết ở nhà, thay vì ở lại công ty để tham dự hay theo dõi trực tiếp các buổi diễn của họ. Không phải tôi không quan tâm, chỉ là\dots\ tôi không giỏi hòa vào bầu không khí ấy.

Thật ra, nếu là những nơi như xe buýt hay tàu điện - những không gian chật hẹp nhưng ai cũng có mục đích riêng, ai làm việc nấy - thì tôi lại chẳng thấy phiền. Nhưng concert thì khác. Ở đó, mọi người hòa mình vào dòng cảm xúc chung, cùng reo hò, cổ vũ, chia sẻ niềm vui với nhau. Còn tôi? Tôi chỉ cảm thấy như một người ngoài cuộc. Tôi đã thử vài lần, nhưng cuối cùng vẫn không thể thay đổi suy nghĩ đó.

Nhấp một ngụm nước, tôi quay trở lại. Và ngay lập tức, đập vào mắt tôi là cảnh sáu người cùng hô lên với một nụ cười rạng rỡ: ``\textbf{Kuon-senpai! Chúc mừng năm mới!}''

``Ối chà, mấy đứa quý hóa ghê nhỉ?'' tôi đáp lại, đặt cốc nước xuống.

Mio-chan hỏi thăm: ``Mấy ngày nay anh ăn Tết với gia đình thế nào?''

``Ờm\dots\ cũng vui.'' Có lẽ tôi không biết phải nhận xét Tết năm nay như thế nào thật.

``Mọi người ở nhà anh vẫn khỏe chứ?'' Matsuri-chan hỏi. Này  này, mấy đứa định hỏi lại đúng những câu mà Tanigo-san vừa hỏi anh đáy à?

Nhưng dù vậy, tôi vẫn giữ một nụ cười, điềm nhiên đáp: ``Ổn. Rất ổn. Cảm ơn.''

Đến lượt Flare-chan, nhưng không phải là hỏi thăm sức khỏe của tôi: ``Em đã được các senpai ở trên kể về trải nghiệm ăn Tết quê anh rồi. Nghe có vẻ vui lắm nhỉ?''

``Anh thấy nó cũng na ná giống Tết Dương lịch ở bên đó thôi. Nhưng, ừm, nó vui.''

``Nó như thế nào vậy?'' Kanata-chan chiêm vào, tỏ vẻ tò mò.

``Nói bây giờ ra hơi khó giải thích, nhưng thôi được, để anh tóm gọn lại nhé\dots''

Thế là tôi bắt đầu thuyết giảng về các phong tục Tết ở quê tôi. Tất cả mọi người, nhất là Marine-chan, Flare-chan và Kanata-chan, đều rất chăm chú lắng nghe. Những người còn lại dù đã trải qua rồi nhưng vẫn không giấu được sự thích thú.

``Giờ em mới biết ở quê anh có cả tập tục thờ cúng tổ tiên nha.'' Thiên sứ-chan nói, gật gù hiểu.

``Đấy là chuyện bình thường mà.''

``Như kiểu, chúc phúc cho những người lên thiên đàng ạ?''

``Nó cũng\dots\ dạng kiểu vậy.'' Tôi không hiểu tại sao em ấy lại dùng cách ví von như vậy.

``Em giờ cũng mới biết cả Ông Công Ông Táo nữa đó,'' Nui-chan lên tiếng, đôi mắt vẫn long lanh, tỏ vẻ ngạc nhiên.

Tôi bật cười: ``Nó nổi tiếng đến mức mà còn tạo nên một chương trình TV cơ. Ví dụ như `Gặp nhau cuối năm'. Chuyên kể những sự kiện nổi bật trong một năm vừa qua. Mỗi tội phải đọc báo nghe tin nhiều mới hiểu được. Chương trình đó còn có mấy bài hát nhép nghe buồn cười lắm.''

Ngay khi tôi dứt lời, Aqua-chan phấn khích nhìn về phía tôi: ``A, nói đến nhạc thì\dots\ Anh có xem buổi biểu diễn của chúng em không vậy!?''

\dots Rồi xong. Chết đứng luôn. Không biết nên trả lời em ấy như thế nào nữa. Biết ngay là sẽ có người hỏi cái này mà.

Nói thật thì sợ mất lòng họ, mà nói dối thì lại bị tra khảo thêm. Mà họ hỏi thế này tức là đã nghi ngờ tôi không xem rồi. Chết tiệt, đúng là khó thoát.

May mắn thay, A-san từ phía sau họ xuất hiện và nói hộ tôi: ``Hikawa-san bận việc gia đình nên cậu ấy không thể xem được.''

Cảm ơn chị nhiều, A-san! Chị đã cứu em một bàn thua trông thấy!

Tất cả các cô gái, đúng như tôi đoán, tỏ vẻ hơi buồn, nhưng họ cũng đồng cảm với tôi.

Chị ấy nói cũng có phần đúng. Mùng 1 Tết cha, mùng 2 Tết mẹ, chưa kể còn đám giỗ của ông nội tôi sắp tới nữa. Tôi chạy cả ngày ngoài đường để thăm họ hàng cũng đủ thấm mệt rồi, lấy đâu ra sức để mà xem nọ xem kia chứ.

A-san nhìn tôi và nói: ``Chúng tôi cũng đã ghi lại toàn bộ buổi concert `Bloom,' ấy rồi, khi nào cậu đi làm trở lại sẽ đưa cho cậu, coi đó như là quà Tết của công ty cho cậu ha.''

\dots Chắc là cũng nên nhận thôi nhỉ? Dù chằng biết tôi có bao giờ xem chứng không. ``Dạ, vâng, em xin nhận.''

``Thế là may rồi nhỉ, Kuon-senpai?'' Akutan nhìn tôi với sự nhẹ nhõm trên mặt. ``Giờ anh không lo sẽ bị bỏ lại phía sau nha!''

``Bọn em có nhiều trò vui trong buổi concert ấy lắm! Khi nào anh về chúng ta cùng xem lại nha!'' Matsuri nói với tông giọng vui vẻ.

Tôi khẽ gật đầu, mỉm cười: ``Được chứ, anh chờ mấy đứa đấy.''

Và rồi, suốt khoảng mười phút sau đó, tất cả mọi người ai nấy đều hoan hỷ, cười nói vui vẻ khi chúng tôi ngồi lại với nhau, kể lại những câu chuyện hài hước mà cả bọn đã từng trải qua.

Một trong số đó lại là một sự hiểu nhầm đi vào lịch sử. ``Thế hóa ra mấy đứa nhầm cái gì cơ? Tuyên ngôn ăn kiêng thành thông báo sa thải á?''

Mio-chan phồng má lên, phụng phịu: ``Không vui tí nào đâu! Tờ giấy đó bị gạch xóa chi chít, mà cách viết của người đó thì\dots\ nó gây lú dã man luôn!''

``Đâu, đưa anh xem nào?''

Matsuri và Marine giơ lên một tờ giấy đầy những nét gạch ngang, gạch chéo, những dòng chữ nắn nót nhưng lại bị tô đè bằng những nét bút nguệch ngoạc. Camera của họ khá nét, nên tôi có thể đọc được rõ từng chữ.

``À, anh biết cái này là ở đâu rồi! Đây đúng là `tuyên ngôn ăn kiêng' của Usake-san!'' Nhìn nét chữ của anh ta là tôi đoán được ngay.

``Thế rốt cuộc là thế nào vậy?'' Flare-chan nhìn lại tờ giấy đó, nghiêng đầu.

``Ừm\dots\ hôm trước, ông ấy than phiền với anh rằng dạo này trông hơi `phệ'. Anh hỏi thử cân nặng thì ông ấy bảo: `Tôi có 90 cân thôi.'

Tôi không nói phét đâu. Tôi đây nặng đến 95 cân đấy, nhưng mà tôi cao, còn ông ấy thì thấp hơn tôi một chút.

Tôi vừa nhớ lại cảnh đó, vừa cười: ``Anh bảo là: `Thế là vừa rồi, cần gì ăn kiêng?' Nhưng ông ấy lại lắc đầu nguầy nguậy: `Không được! Kỳ này tôi sẽ nhất quyết ăn kiêng một phen!' Và thế là\dots\ cái này ra đời.''

``Thế rồi kết quả thế nào?'' Marine hỏi, bụm miệng cười.

``Nhịn được đúng ba ngày. Đến ngày thứ tư thì kêu đói quá chịu không nổi nữa, thế là\dots\ bỏ!''

Tất cả sáu cô gái gần như đồng loạt ngửa người ra sau, bật cười nghiêng ngả. ``Anh nghĩ ông ấy gạch chi chít như thế vì sợ bị quê, nhưng mà ông ấy có thể ném nó vào máy hủy giấy được mà!''

Akutan cười trong nước mắt: ``Trời ơi, thế mà cũng bày đặt ăn kiêng nữa!''

Matsuri-chan thì vừa cười vừa đập bàn: ``Không có chí tiến thủ gì hết!''

``Thì, Matsuri hôm đó cũng chơi bài nhịn ăn như thế năm ngoái\footnote{Xem lại {\flaregothic ÉPISODE 33}, quyển số Ba.} tới mức suýt ngỏm đấy. Nhớ hồi đó không?''

Thấy rằng mình bị gọi tên, em ấy cũng chỉ cười trừ. ``T-Tại hồi đó em không biết ăn kiêng thế nào\dots''

``Mà, lúc đó mấy đứa định làm gì với tờ giấy này?'' tôi hỏi, tỏ vẻ hiếu kỳ với tờ tuyên ngôn đó. ``Anh mở lại máy hôm qua thì thấy Flare-chan gọi điện với gửi tin cho anh này.''

``Thì thế! Nhưng mà Mio-senpai bảo anh không nhấc máy, cũng chẳng trả lời tin nhắn. Sao lại thế?'' nàng Dị giới hỏi, tỏ vẻ hơi bất bình dù môi em ấy vẫn đang cười trừ.

``À, đúng rồi! Nghỉ Tết thì công việc phải vứt sang một bên chứ! Anh chỉ giữ liên lạc với A-san và Tanigo-san thôi. Mà kể cả thế, anh cũng dặn họ trước rồi: chỉ khi nào họp toàn công ty hoặc có việc đột xuất cần đến anh thì hẵng gọi, chứ bình thường thì không.'' Đó đã là một trong những quy tắc tôi đã đặt ra ngay từ năm nay kể từ sau sự kiện bất ngờ đó.

``Nhưng nếu là việc của tụi em thì sao? Em tưởng là anh đã bỏ rơi bọn em rồi đấy\dots'' Kanata-chan hỏi, tỏ vẻ bất mãn.

Tôi nhún vai đáp lại: ``Thì tự giải quyết với nhau. Anh đâu có lúc nào cũng kè kè bên các em để xử lý mọi chuyện đâu? Mấy đứa cũng phải trưởng thành lên đi.''

``Nhưng mà sau đó họ cứ đè đầu cưỡi cổ em này!'' em ấy rít lên, hoàn toàn bất mãn. Bộ em ấy có gì đó uất ức à?

``Sao thế? Kể anh nghe xem.''

Thiên sứ liền bức xúc kể lại việc bị lôi ra làm ``vật thí nghiệm'' cho đủ thứ ý tưởng tập luyện thần tượng kỳ quái, từ nhảy dây liên tục trong khi hát đến việc tập múa với tạ đeo ở chân. Thậm chí, chiếc vầng hào quang của em ấy còn bị lôi ra làm đồ chơi.

Tôi bật cười: ``Trời ơi, tưởng gì ghê gớm lắm! Anh thấy mấy đứa toàn làm mấy trò này suốt mà. Đúng là chuẩn hội hề của công ty luôn!''

Sau khi cả nhóm cười rôm rả về chuyện ``tuyên ngôn ăn kiêng,'' bầu không khí dần lắng xuống. Nhưng vẫn còn đó những điều chưa nói hết. Tôi nhận ra Kana-chan đang bĩu môi, còn Aqua-chan thì khoanh tay đầy tự hào, như thể em ấy đang có một điều gì đó rất quan trọng cần chia sẻ.

``Em cũng đã nói với mọi người thế nào là idol rồi!'' nàng Hầu gái nói với sự hãnh diện.

``Thế nó là thế nào?''

``Thiên tài tổng vệ sinh nhưng vụng về''-chan nhấp một ngụm nước, hắng giọng đầy khí thế trước khi lên tiếng. ``Một idol mà em hướng đến là người có thể làm những gì mình thích, và quan trọng hơn cả là có thể cống hiến cho khán giả để làm họ vui. Ít nhất, đó là quan điểm của em. Ai cũng có thể có cách xây dựng idol của riêng mình!''

Giọng Aqua-chan chắc nịch, không chút do dự. Tôi khẽ gật đầu, cảm nhận được sự nghiêm túc trong lời nói của em ấy.

``Đúng rồi, chuẩn đấy. Idol có nhất thiết phải là người đứng trên sân khấu, tỏa sáng rực rỡ trước hàng ngàn ánh đèn không? Aqua-chan nói đúng rồi đó.''

Mat-chan khoanh tay, nhìn tôi với vẻ nửa đùa nửa thật. ``Đến cả Kuon-senpai cũng nói đúng nữa kìa. Hiếm có nha!''

Tôi cười khẽ. Cái phong cách nói chuyện kiểu ``cà khịa'' này đúng là không bao giờ thay đổi. Nhưng lần này tôi không để em ấy trêu chọc mà tiếp tục chủ đề.

``Nhưng anh cũng phải nói thêm, mà chắc mấy đứa cũng biết rồi, rằng cách trở thành idol thì tùy mỗi người, nhưng trách nhiệm của nó dù theo cách nào thì rất là cao đấy.''

``Không, cái đó thì bọn em hiểu,'' Mio-chan nói với giọng chắc chắn, đại diện cho cả nhóm khẳng định điều này.

Tôi nhìn quanh một lượt, thấy ai cũng có vẻ nghiêm túc hơn, bèn tiếp tục: ``Theo anh thấy, idol có nghĩa là thần tượng, tức là một hình mẫu mà người khác muốn noi theo, một con người gương mẫu mà mọi người có thể học hỏi. Nhưng idol không nhất thiết phải là người nổi tiếng trên mạng hay trên sân khấu. Idol - hay thần tượng có thể đơn giản là một người trong gia đình, một người bạn nào đó, hay thậm chí là thầy cô giáo.''

Tôi ngừng lại một chút để họ có thời gian suy nghĩ. Dù là những điều đơn giản, nhưng đôi khi phải nghe ai đó nói ra thì mới thật sự ngẫm được.

Marine-chan là người đầu tiên lên tiếng: ``Vậy\dots\ ý của anh là\dots?''

``Ý anh là anh không quan tâm các em xây dựng hình tượng idol như thế nào. Đó là lựa chọn của các em. Nhưng điều anh muốn là các em phải xây dựng nó sao cho phù hợp với cá tính của chính mình, đồng thời vẫn phải văn minh, có trách nhiệm và phù hợp với quy tắc đạo đức xã hội. Hiểu chửa?''

Tất cả mọi người đều gật đầu. Không cần nói thêm gì, tôi biết họ đã hiểu được thông điệp mà mình muốn truyền tải.

Tôi chậm rãi dựa lưng vào ghế, nhìn vào màn hình nơi những gương mặt quen thuộc của \hll\ đang hiện lên. ``Idol'' không chỉ đơn thuần là một nghề nghiệp hay một danh hiệu, mà còn là một lý tưởng, một trách nhiệm mà họ đã lựa chọn gánh vác. Và tôi hy vọng, những lời nói này sẽ giúp họ vững bước hơn trên con đường ấy.

``Idol là một điều gì đó rất rộng lớn, rất tự do, rất vui vẻ\dots\ nhưng đồng thời cũng đầy áp lực. Cứ chờ đến khi trưởng thành thêm chút nữa, các em sẽ hiểu thôi.''

Và đó chính là những bài học về idol mà tôi muốn truyền tải cho các thành viên nhân dịp Tết Tân Sửu này.

Vẫn còn rất nhiều câu chuyện nữa mà chúng tôi có thể kể, nhưng có vẻ như cuộc trò chuyện hôm nay phải tạm dừng tại đây\dots

\dots nếu không vì một ai đó mà đáng ra không nên xuất hiện trước mặt họ phá đám.

``\textbf{Kuon-kun\~{}!! Bánh quy của em làm xong rồi đây!! Anh có muốn ăn thử không!?}'' Một giọng nói lanh lảnh vang lên từ phía cầu thang.

Tôi chết sững. Chết mợ, là Aloe-chan!

Cô nàng succubus bất ngờ xuất hiện, bưng một khay bánh quy vừa làm từ trên bếp xuống, trông đầy tự hào với thành phẩm của mình. Ngay khi giọng nói của em ấy vang lên, bầu không khí trong cuộc gọi bỗng chốc thay đổi. Từ phía màn hình, tôi có thể thấy những ánh mắt nghi hoặc bắt đầu lóe lên.

Nàng Succubus nghiêng đầu nhìn màn hình máy tính, tỏ vẻ tò mò: ``Anh đang nói chuyện với ai vậy, Kuon-kun? Để em xem thử-''

Tôi vội lao tới chắn em ấy, toát mồ hôi hột: ``K-Không có gì đâu!''

Ở phía bên kia đầu dây, Kana-chan nhướng mày: ``Ủa? Ai vậy, Kuon-senpai?''

Akutan thì híp mắt đầy nghi ngoặc: ``Anh cũng có bạn gái rồi à?''

Tôi lắp bắp: ``Ờm\dots\ Không\dots? Bọn anh chỉ là bạn bè thôi.''

Hẳn nhiên, câu nói đó không thuyết phục được họ. Marine-chan nghe vậy càng muốn tỏ vẻ trêu chọc tôi nhiều hơn: ``Anh đã có bạn gái rồi ư\~{}?''

Tôi liền bật dậy, vừa cố gắng chắn Aloe-chan vừa nói: ``Không có! Mấy đứa nghĩ vớ vẩn gì thế!''

Cả Mio-chan và Mat-chan cũng không đứng về phía tôi. Họ chỉ đồng thanh nhìn tôi và bóng dáng ở phía sau: ``Nghi lắm nha\dots''

Nhìn thấy ánh mắt tò mò của bọn họ, tôi biết nếu còn tiếp tục thì chắc chắn mình sẽ bị tra hỏi đến cùng.

``Thôi được, mấy bữa nữa anh đi làm thì ta kể tiếp nha! Bảo trọng!'' Nói rồi, tôi lập tức kết thúc cuộc trò chuyện mà không để cho ai kịp phản ứng gì.

Trời ạ, Aloe-chan à\dots\ Em chưa sẵn sàng để có thể gặp mặt bọn họ được đâu\dots

Và thế là, cuộc gọi kết thúc. Một cách rất đường đột.

Tôi chớp mắt, cảm thấy bầu không khí xung quanh bỗng trở nên im lặng đến lạ thường. Bên đầu dây bên kia, chắc hẳn mọi người đang vô cùng khó hiểu. Mà không chỉ họ, ngay cả tôi cũng thế. Tự dưng lại cúp máy như vậy là sao chứ?

Vài giây trôi qua, tôi thở dài, quay sang Aloe-chan - người vẫn đang đứng đó, chớp mắt ngây thơ như chưa hiểu chuyện gì. Tôi thở dài một tiếng: ``\dots Em có thể đừng hét to thế không?''

Cô nàng nhíu mày, nghiêng đầu đầy vô tội: ``Ơ? Nhưng em đã làm bánh quy cho anh mà! Em muốn được anh nghiệm thử thành quả do chính tay em làm!''

Ngay lúc đó, một ánh mắt sắc bén khác quét qua tôi. Ryuuhou từ nãy giờ đứng khoanh tay quan sát, cuối cùng cũng lên tiếng: ``\dots Anh vừa gọi cho ai thế?''

Tôi im lặng trong thoáng chốc, lướt mắt sang Aloe-chan. Tôi sẽ không nói chừng nào em ấy vẫn còn ở đây.

Nhìn từ góc độ của em ấy, chắc chắn em ấy đang nghĩ tôi có một bí mật gì đó mà em không có quyền được biết.

Nhưng thay vì tiếp tục để bị tra hỏi, tôi thản nhiên cầm một chiếc bánh quy lên, cắn một miếng, rồi gật đầu giả ngơ: ``Không có gì, ăn bánh thôi.''

Nàng Succubus còn chưa kịp phản ứng thì Ryuuhou cũng đã với lấy một chiếc bánh, nhấm nháp thử. ``Hừm\dots\ cũng ngon đấy.''

Đôi mắt cô lập tức sáng rực với lời khen bất ngờ đó: ``Thật á!?''

Và, chỉ với một câu nhận xét ngắn gọn đó, mọi chuyện vừa xảy ra bị quẳng ra sau đầu. Aloe-chan nhanh chóng ngồi xuống cùng chúng tôi, hào hứng thưởng thức thành phẩm do chính tay mình làm.

Tôi liếc nhìn em ấy một chút rồi lặng lẽ nhâm nhi bánh quy, uống trà như thể chẳng có chuyện gì đặc biệt vừa diễn ra. Hy vọng là sáu người đó không để ý tới em ấy.

Cầu trời khấn phật.

\dots Được rồi, tôi phải thừa nhận là chuyện vừa rồi khá là hỗn loạn. Suýt chút nữa là lộ hết bí mật của tôi với Tanigo-san rồi.

Mà khoan, lúc nãy tôi đang định nói cái gì để kết thúc câu chuyện này nhỉ? À, phải\dots *e hèm*

Mùa đông cũng sắp qua đi. Một mùa xuân mới lại sắp đến, mang theo không khí nhộn nhịp và biết bao điều mới mẻ. Cũng đồng nghĩa với việc sẽ lại có thêm nhiều công việc nữa đang chờ tôi tại chính \hll\ này.

Sau \textit{Bloom,} sẽ là nhiều buổi concert khác nữa. Sẽ có thêm nhiều thành viên mới gia nhập, sẽ có thêm nhiều câu chuyện vui, cũng như những khoảnh khắc buồn đáng nhớ.

Nhưng dù thế nào đi nữa, tôi vẫn là một phần của đại gia đình \hll\ - một nơi tràn ngập sự tréo ngoe, lắm chuyện điên rồ, nhưng cũng đầy ắp sự nồng nhiệt và tiếng cười.

Và tôi sẽ tiếp tục cố gắng hơn nữa!

Chúc mừng năm mới Tân Sửu!

\pov{-}

Bên này, sau khi màn hình cuộc gọi vụt tắt đầy đột ngột, sáu cô gái vẫn còn sững sờ.

Matsuri chớp mắt mấy lần, rồi búng tay đánh ``tách'' một cái: ``Thế là sao? Chả phải Kuon-senpai vừa bị ai đó bắt quả tang à?''

``Ờ đúng, nhưng mà là ai nhỉ?'' Aqua nghiêng đầu, tay chống cằm ra vẻ suy luận. ``Nghe giọng thì\dots\ trẻ, nhưng mà có gì đó hơi-''

``Bánh quy,'' Mio đột nhiên lên tiếng, khiến mọi người quay sang nhìn. Cô gật gù như đã hiểu ra điều gì. ``Cô bé đó nói gì nhỉ? `Bánh quy của em làm xong rồi đây!!' hay là, kiểu như thế sao?''

``Chậc, người quen của Kuon-senpai à?'' Marine huých nhẹ vào Flare, nhưng Flare chỉ nhún vai. ``Không nhớ ra là ai, nhưng giọng nghe quen lắm\dots''

Kanata chống nạnh, đôi cánh nhỏ rung rung đầy suy nghĩ. ``Một cô bé có giọng khá trẻ con, năng động, biết làm bánh quy\dots\ hừm\dots''

Mọi người nhìn nhau, cố lục lại trí nhớ. Cái tên kia\dots\ cái gì mà ``Ma'' ấy nhỉ? Hay là ``A'' nhỉ?

``\dots Mà khoan, anh ấy hoảng hốt lắm đúng không?'' nàng Lễ hội đột nhiên reo lên, đôi mắt lấp lánh. ``Chị thề là anh ấy chưa bao giờ phản ứng kiểu đó đâu nha! Chắc chắn là có chuyện mờ ám!''

Nàng Hầu gái gật đầu liên tục. ``Chắc luôn! Biểu cảm của Kuon-senpai lúc đó kiểu `Thôi xong!', trông rõ ràng quá mà!''

Trong khi nhóm các cô gái vẫn đang bàn tán rôm rả, cách đó không xa, A-chan thở dài, đẩy gọng kính lên. Cô quay sang người đàn ông đứng bên cạnh, người vẫn giữ nguyên nụ cười có phần bất lực trên gương mặt.

``YAGOO\dots\ kế hoạch của ông suýt bị bại lộ rồi đấy.''

Tanigo chỉ cười trừ, khoanh tay nhìn màn hình cuộc gọi giờ đã trống không: ``Nhưng cuối cùng thì đâu có bại lộ, đúng không?''

A-chan lắc đầu, nhưng cũng chẳng nói thêm gì. Dù sao thì, đây cũng chẳng phải lần đầu tiên có chuyện như vậy xảy ra\dots

Cuối cùng, Kuon đã phải đứng ra giải thích hết cho mọi chuyện, nói dối rằng đó là một cô em họ của cậu sang chơi, tới lúc đó mọi người mới thôi không tra hỏi thêm nữa, dù rằng họ vẫn còn rất nhiều câu hỏi trong đầu.

Không biết liệu năm Tân Sửu này Hikawa Kuon còn có những tam tai nào nữa đây?

\end{document}